\chapter{Adams differentials}
\label{ch:Adams}

The goal of this chapter is to describe the values of the
Adams differentials in the motivic Adams spectral sequence.
These values are given in Tables \ref{tab:Adams-d2}, \ref{tab:Adams-d3},
\ref{tab:Adams-d4}, \ref{tab:Adams-d5}, and \ref{tab:Adams-higher}.
See also the Adams charts in \cite{Isaksen14a} for a graphical
representation of the computations.
\revv{
For easy reference, the many lemmas in this chapter are labelled with degrees that match the degrees given in the tables.}

\section{The Adams $d_2$ differential}
\label{sctn:Adams-d2}

\revv{
Table \ref{tab:Adams-d2} lists all of the multiplicative generators 
of the Adams $E_2$-page through the 95-stem.  The third column indicates
the value of the $d_2$ differential, if it is non-zero.
A blank entry in the third column indicates that the $d_2$ differential
is zero.
The fourth column indicates the proof. A blank entry in the fourth column
indicates that there are no possible values for the differential.
The fifth column gives alternative names for the element, as used 
in \cite{Tangora70a}, \cite{Bruner97}, and \cite{Isaksen14c}.
}

\begin{thm}
\label{thm:Adams-d2}
Table \ref{tab:Adams-d2} lists the values of the
Adams $d_2$ differential on all multiplicative generators
through the 95-stem.
\end{thm}

\rev{
\begin{remark}
\label{rem:d2}
A previous version of this manuscript left uncertain the value of
$d_2$ on three multiplicative generators.  These three values
have since been determined by Dexter Chua \cite{Chua21}.
We have included those values here, but we defer to \cite{Chua21}
and \cite{BF21} for their proofs.
Also, we are grateful to
Joey Beauvais-Feisthauer \cite{BF21} for discovering an error
in our previous calculation of $d_2(x_{85,6})$.
\end{remark}
}

\begin{proof}
The fourth column of Table \ref{tab:Adams-d2} gives information on the
proof of each differential.
Most follow immediately by comparison to
the Adams spectral sequence for $C\tau$
\revv{\cite{IWX19}}.
A few additional differentials follow by comparison to the classical
Adams spectral sequence for $\tmf$ \revv{\cite{BR21}}.

If an element is listed in the fourth column of Table \ref{tab:Adams-d2}, 
then the corresponding differential can be deduced from a straightforward
argument using a multiplicative relation.  For example, it is possible
that $d_2(\D h_1 h_3)$ equals $\tau d_0 e_0$.  However,
$h_0 \cdot \D h_1 h_3$ is zero, while $h_0 \cdot \tau d_0 e_0$ is non-zero.
Therefore, $d_2(\D h_1 h_3)$ must equal zero.

In some cases, it is necessary to combine these different techniques to
establish the differential.

The remaining more difficult computations are carried out in the
following lemmas.  We refer to \cite{Chua21} and \cite{BF21} in
a few cases.
\end{proof}


\begin{lemma}
\label{lem:d2-Dx}
\revdeg{61, 9, 32}
$d_2(\D x) = h_0^2 B_4 + \tau M h_1 d_0$.
\end{lemma}

\begin{proof}
We have a differential $d_2(\D x) = h_0^2 B_4$ in the
Adams spectral sequence for $C\tau$.
Therefore, $d_2(\D x)$ equals either
$h_0^2 B_4$ or $h_0^2 B_4 + \tau M h_1 d_0$.

We have the relation
$h_1^2 \cdot \D x = P h_1 \cdot \tau \D_1 h_1^2$
\revv{(as usual, we rely on complete information about classical products in a large range \cite{Bruner97} \cite{BR20})}, so
$h_1^2 d_2(\D x) = P h_1 d_2(\tau \D_1 h_1^2) = 
P h_1 h_3 \cdot M h_1 = \tau M h_1^3 d_0$.
Therefore, $d_2(\D x)$ must equal
$h_0^2 B_4 + \tau M h_1 d_0$.
\end{proof}

\begin{remark}
The proof of \cite{Isaksen14c}*{Lemma 3.50}
is incorrect.  We claimed that
$h_1^2 \cdot \D x$ equals $h_3 \cdot \D^2 h_1 h_3$,
when in fact
$h_1^2 \cdot \D x$ equals $\tau h_3 \cdot \D^2 h_1 h_3$.
\end{remark}

\begin{lemma}
\label{lem:d2-x77,7}
\revdeg{77, 7, 40}
$d_2(x_{77,7}) = \tau M h_1 h_4^2$.
\end{lemma}

The following proof was suggested to us by Dexter Chua.

\begin{proof}
This follows from the interaction between algebraic squaring
operations and classical Adams differentials \cite{Bruner84}*{Theorem 2.2},
applied to the element $x$ in the $37$-stem.
The theorem says that
\[
d_* \Sq^2 x = \Sq^3 d_2 x \dotplus h_0 \Sq^3 x.
\]
The notation means that there is an Adams differential on $\Sq^2 x$
hitting either $\Sq^3 d_2 x = 0$ or $h_0 \Sq^3 x$,
depending on which element has lower Adams filtration.
Therefore
$d_2 \Sq^2 x = h_0 \Sq^3 x$.

Next, observe from \cite{Bruner04} that
$\Sq^3 x = h_0^2 x_{76,6} + \tau^2 d_1 g_2$, so
\[
h_0 \Sq^3 x = h_0^3 x_{76,6} = \tau M h_1 h_4^2.
\]
Therefore, there is a $d_2$ differential whose value is
$\tau M h_1 h_4^2$, and the possibility is that
$d_2(x_{77,7})$ equals $\tau M h_1 h_4^2$.
\end{proof}

\begin{lemma}
\label{lem:d2-tB5g}
\revdeg{86, 14, 47}
$d_2(\tau B_5 g) = \tau M h_0^2 g^2$.
\end{lemma}

\begin{proof}
We use the Mahowald operator methods of 
Section \ref{subsctn:Mahowald-operator}.
\revv{According to \cite{Isaksen19}*{Table 1}},
the map $p_*: \Ext_\C \map \Ext_B$ takes
$d_0 \cdot \tau B_5 g$ to
$\tau h_0 a g^3 v_3^2$, which is non-zero.
We deduce that the product $d_0 \cdot \tau B_5 g$ is non-zero
in $\Ext$.
\revv{By inspection of motivic weights},
the only possibility is that it equals $\tau M g \cdot h_0 m$.

Now $d_2(\tau M g \cdot h_0 m)$ equals
$\tau M g \cdot h_0^2 e_0^2$, which we also know is non-zero
since it maps to the non-zero element
$\tau h_0^2 d e g^2 v_3^2$ of $\Ext_B$
\revv{by \cite{Isaksen19}*{Theorem 1}}.
It follows that $d_2(\tau B_5 g)$ 
\revv{is non-zero.  By inspection of motivic weights, the only
possibility is $\tau M h_0^2 g^2$.}
\end{proof}

\section{The Adams $d_3$ differential}
\label{sctn:Adams-d3}

\revv{
Table \ref{tab:Adams-d3} lists the multiplicative generators 
of the Adams $E_3$-page through the 95-stem
whose $d_3$ differentials are non-zero, or whose $d_3$ differentials
are zero for non-obvious reasons.
}

\begin{thm}
\label{thm:Adams-d3}
Table \ref{tab:Adams-d3} lists some values of the
Adams $d_3$ differential on multiplicative generators.
Through the 95-stem, 
the Adams $d_3$ differential is zero on all multiplicative
generators not listed in the table.
\end{thm}

\rev{
\begin{remark}
A previous version of this manuscript left uncertain the value of
$d_3$ on several multiplicative generators.  These values
have since been determined by Dexter Chua \cite{Chua21}.
We have included those values here, but we defer to \cite{Chua21}
for their proofs.
We are also grateful to Dexter Chua for correcting a few mistakes
in the values of the $d_3$ differential.
\end{remark}
}

\begin{proof}
The $d_3$ differential on many multiplicative generators is zero.
A few of these multiplicative generators appear
in Table \ref{tab:Adams-d3} because their proofs require further
explanation.
For the remaining majority of such multiplicative generators,
the $d_3$ differential is zero
because there are no possible non-zero
values, because of comparison to the Adams spectral sequence for $C\tau$,
or because the element is already known to be a permanent cycle
as shown in Table \ref{tab:Adams-perm}.
These cases do not appear in Table \ref{tab:Adams-d3}.

The last column of Table \ref{tab:Adams-d3} gives information on the
proof of each differential.
Most follow immediately by comparison to
the Adams spectral sequence for $C\tau$.  
A few additional differentials follow by comparison to the classical
Adams spectral sequence for $\tmf$, or by comparison to the
$\C$-motivic Adams spectral sequence for $\mmf$.

If an element is listed in the last column of Table \ref{tab:Adams-d3}, 
then the corresponding differential can be deduced from a straightforward
argument using a multiplicative relation.  For example, 
\[
d_3( h_1 \cdot \tau P d_0 e_0 ) =
P h_1 \cdot d_3(\tau d_0 e_0) =
P^2 h_1 c_0 d_0,
\]
so $d_3(\tau P d_0 e_0)$ must equal $P^2 c_0 d_0$.

If a $d_4$ differential is listed in the last column of 
Table \ref{tab:Adams-d3}, then
the corresponding differential is forced by consistency with
that later differential.
In each case,
a $d_3$ differential on an element $x$
is forced by the existence of a later
$d_4$ differential on $\tau x$.
For example, Table \ref{tab:Adams-d4} shows that
there is a differential
$d_4(\tau^2 e_0 g) = P d_0^2$.
Therefore, $\tau e_0 g$ cannot survive to the $E_4$-page.
It follows that $d_3(\tau e_0 g) = c_0 d_0^2$.

In some cases, it is necessary to combine these different techniques to
establish the differential.

The remaining more difficult computations are carried out in the
following lemmas.  We refer to \cite{Chua21} in
a few cases.
\end{proof}

\begin{prop}
\label{prop:Adams-perm}
Some permanent cycles in the
$\C$-motivic Adams spectral sequence are shown in 
Table \ref{tab:Adams-perm}.
\end{prop}

\begin{proof}
The third column of the table gives information on the proof
for each element.  If a Toda bracket is given in the third
column, then
the Moss Convergence Theorem \ref{thm:Moss} implies that the
element must survive to detect that Toda bracket
(see Table \ref{tab:Toda} for more information on how each
Toda bracket is computed).
If a product is given in the third column, then
the element must survive to detect that product
(see Table \ref{tab:misc-extn} for more information
on how each product is computed).
In a few cases, the third column refers to a specific lemma
that gives a more detailed argument.
\end{proof}

\begin{lemma}
\label{lem:d3-h2h5}
\mbox{}
\begin{enumerate}
\item
\revdeg{34, 2, 18}
$d_3(h_2 h_5) = \tau h_1 d_1$.
\item
\revdeg{74, 6, 38}
$d_3(P h_2 h_6) = \tau h_1 h_4 Q_2$.
\end{enumerate}
\end{lemma}

\begin{proof}
In the Adams spectral sequence for $C\tau$, there is an
$\eta$ extension from $h_2 h_5$ to $\ol{h_1^2 d_1}$.
The element $\ol{h_1^2 d_1}$ maps to $h_1^2 d_1$ under projection
from $C\tau$ to the top cell, so
$h_2 h_5$ must also map non-trivially under projection from
$C\tau$ to the top cell.
The only possibility is that $h_2 h_5$ maps to $h_1 d_1$.
Therefore, $\tau h_1 d_1$ must be hit by a differential.
This establishes the first differential.

The proof for the second differential is identical, using that
there is an $\eta$ extension from $P h_2 h_6$ to $\ol{h_1^2 h_4 Q_2}$
in the Adams spectral sequence for $C\tau$.
\end{proof}

\begin{lemma}
\label{lem:d3-t^2D1h1^2}
\revdeg{54, 6, 28}
$d_3(\tau^2 \D_1 h_1^2) = \tau M c_0$.
\end{lemma}

\begin{proof}
The element $M P$ maps to zero under inclusion of the bottom cell into 
$C \tau$. 
Therefore,
$M P$ is either hit by a differential, or it is the target of a 
hidden $\tau$ extension.  
\revv{If it is the target of a hidden $\tau$ extension},
then the only possibility is that $\tau M c_0$ is zero in the
$E_\infty$-page, and that there is a hidden
$\tau$ extension from $M c_0$ to $M P$.

\revv{It remains to show that $M P$ cannot be hit by a differential.
The only possibility is that $d_4(\tau^2 \D_1 h_1^2)$ might
equal $M P$. 
Note that $P h_1 \cdot \tau^2 \D_1 h_1^2$ 
equals $h_1 ( \tau h_1 \cdot \Delta x)$ in the $E_4$-page.
This means that 
$d_4(P h_1 \cdot \tau^2 \D_1 h_1^2)$ cannot equal $M P^2 h_1$
since $M P^2 h_1$ is not divisible by $h_1$.
In turn, $d_4(\tau^2 \D_1 h_1^2)$ cannot equal $M P$.
}
\end{proof}

\begin{lemma}
\label{lem:d3-tDg2.h0^3}
\revdeg{68, 11, 35}
$d_3(\tau h_0^3 \cdot \D g_2) = \tau^3 \D h_2^2 e_0 g$.
\end{lemma}

\begin{proof}
Table \ref{tab:unit-mmf} shows that the element
$h_0 h_5 i$ maps to $\D^2 h_2^2$ in
the Adams spectral sequence for $\tmf$.

Now $\D^2 h_2^2 d_0$ is not zero and not divisible by $2$ in $\tmf$.
Therefore, $\kappa \{ h_0 h_5 i\}$ must be non-zero and 
not divisible by $2$ in $\pi_{68,36}$.
The only possibility is that
$\kappa \{h_0 h_5 i\}$ is detected by 
$P h_2 h_5 j = d_0 \cdot h_0 h_5 i$, and that
$P h_2 h_5 j$ is not an $h_0$ multiple in the $E_\infty$-page.
Therefore, $\tau \D g_2 \cdot h_0^3$ cannot survive to the
$E_\infty$-page.
\end{proof}

\rev{
\begin{lemma}
\label{lem:d3-tD3'}
\revdeg{69, 8, 36}
$d_3(\tau D'_3) = \tau^2 M h_2 g$.
\end{lemma}
}

\begin{proof}
Table \ref{tab:Toda} shows that
the Toda bracket $\langle 2, 8 \sigma, 2, \sigma^2 \rangle$
contains $\tau \nu \kappabar$, which is detected by $\tau^2 h_2 g$.
Table \ref{tab:nu-extn} shows that
$M h_2$ detects $\nu \alpha$ for some
$\alpha$ in $\pi_{45,24}$ detected by $h_3^2 h_5$.
(Beware that there is a crossing extension, $M h_2$ does not
detect $\nu \alpha$ for every $\alpha$ that is detected by
$h_3^2 h_5$.)
It follows that
$\tau^2 M h_2 g$ detects
$\langle 2, 8 \sigma, 2, \sigma^2 \rangle \alpha$.

This expression is contained in
$\langle 2, 8 \sigma, \langle 2, \sigma^2, \alpha \rangle \rangle$.
Lemma \ref{lem:2,sigma^2,theta4.5} shows that
the inner bracket equals
$\{0, 2 \tau \kappabar^3 \}$.

The Toda bracket
$\langle 2, 8 \sigma, 0 \rangle$ 
in $\pi_{68,36}$
consists entirely of multiples of $2$.
The Toda bracket
$\langle 2, 8 \sigma, 2 \tau \kappabar^3 \rangle$ contains
$\langle 2, 8 \sigma, 2 \rangle \tau \kappabar^3$.
This last expression equals zero because 
\[
\langle 2, 8 \sigma, 2 \rangle = \tau \eta \cdot 8 \sigma = 0
\]
by Corollary \ref{cor:2-symmetric}.
Therefore, 
$\langle 2, 8 \sigma, 2 \tau \kappabar^3 \rangle$ equals
its indeterminacy, which
consists entirely of multiples of $2$ in $\pi_{68,36}$.

We conclude that $\tau^2 M h_2 g$ is either hit by a differential, or
is the target of a hidden $2$ extension.  
Lemma \ref{lem:2-h3A'} shows that
there is no hidden $2$ extension from $h_3 A'$ to $\tau^2 M h_2 g$,
and there are no other possible extensions to $\tau^2 M h_2 g$.

Therefore, $\tau^2 M h_2 g$ must be hit by 
a differential, and the only possible source of this differential
is $\tau D'_3$.
\end{proof}

\begin{lemma}
\label{lem:d3-t^2Mh0l}
\revdeg{77, 14, 40}
$d_3(\tau^2 M h_0 l) = \D^2 h_0 d_0^2$.
\end{lemma}

\begin{proof}
Table \ref{tab:Adams-d4} shows that
$d_4(\tau^2 d_0 e_0 + h_0^7 h_5)$ equals $P^2 d_0$, so
$d_4(\tau^3 M h_1 d_0 e_0)$ equals $\tau M P^2 h_1 d_0$.
We have the relation $h_0 \cdot \tau^2 M h_0 l = \tau^3 M h_1 d_0 e_0$,
but the element $\tau M P^2 h_1 d_0$ is not divisible by $h_0$.
Therefore, $\tau^2 M h_0 l$ cannot survive to the $E_4$-page.

By comparison to the Adams spectral sequence for $\tmf$,
the value of $d_3(\tau^2 M h_0 l)$ cannot be
$\tau^3 \D h_1 e_0^3 + \D^2 h_0 d_0^2$ or
$\tau^3 \D h_1 e_0^3$.
The only remaining possibility is that
$d_3(\tau^2 M h_0 l)$ equals
$\D^2 h_0 d_0^2$.
\end{proof}

\begin{lemma}
\label{lem:d3-h0^3x78,10}
\revdeg{78, 13, 40}
$d_3(h_0^3 x_{78,10}) = \tau^6 e_0 g^3$.
\end{lemma}

\begin{proof}
Suppose that $h_0^3 x_{78,10}$ were a permanent cycle.
Then it would map under inclusion of the bottom cell
to the element $h_0^3 x_{78,10}$ in the 
Adams $E_\infty$-page for $C\tau$.

There is a hidden $\nu$ extension from $h_0^3 x_{78,10}$
to $\D^3 h_1^2 h_3$ in the Adams $E_\infty$-page for $C\tau$.
Then $\D^3 h_1^2 h_3$ would also have to be in the image
of inclusion of the bottom cell.  
The only possible pre-image is the element $\D^3 h_1^2 h_3$
in the Adams spectral sequence for the sphere, but this element
does not survive by Lemma \ref{lem:d4-D^3h1^2h3}.

By contradiction, we have shown that $h_0^3 x_{78,10}$
must support a differential.  The only possibility is that
$d_3(h_0^3 x_{78,10})$ equals $\tau^6 e_0 g^3$.
\end{proof}

\begin{lemma}
\label{lem:d3-x1}
\revdeg{79, 5, 42}
$d_3(x_1) = \tau h_1 m_1$.
\end{lemma}

\begin{proof}
This follows from the interaction between algebraic squaring
operations and classical Adams differentials \cite{Bruner84}*{Theorem 2.2}.
The theorem says that
\[
d_* \Sq^1 e_1 = \Sq^3 d_3 e_1 \dotplus h_1 \Sq^3 e_1.
\]
The notation means that there is an Adams differential on $\Sq^1 e_1$
hitting either $\Sq^3 d_3 e_1$ or $h_1 \Sq^3 e_1$,
depending on which element has lower Adams filtration.
Therefore
$d_3 \Sq^1 e_1 = h_1 \Sq^3 e_1$.

Finally, we observe from \cite{Bruner04} that
$\Sq^1 e_1 = x_1$ and $\Sq^3 e_1 = m_1$.
\end{proof}

\begin{lemma}
\label{lem:perm-D^2d1}
\revdeg{80, 12, 42}
The element $\D^2 d_1$ is a permanent cycle.
\end{lemma}

\begin{proof}
The element $\D^2 d_1$ in the Adams $E_\infty$-page for $C\tau$
must map to zero under the projection from $C\tau$ to the top cell.
The only possible value in sufficiently high filtration is
$\tau^2 \D h_1 e_0^2 g$.  However, comparison to $\mmf$
shows that this element is not annihilated by $\tau$, and therefore
cannot be in the image of projection to the top cell. 

Therefore, $\D^2 d_1$ 
must be in the image of the inclusion of the bottom cell into $C\tau$.
The element $\D^2 d_1$ is the 
only possible pre-image in the Adams $E_\infty$-page for the sphere 
in sufficiently low filtration.
\end{proof}

\rev{
\begin{lemma}
\label{lem:d3-D^3h1h3}
\revdeg{80, 14, 41}
$d_3(\D^3 h_1 h_3)$ equals either
$\tau^4 \D h_1 e_0^2 g$, $\tau \D^2 h_0 d_0 e_0$, or
$\tau^4 \D h_1 e_0^2 g + \tau \D^2 h_0 d_0 e_0$; and it is not equal to
$d_3(\tau^2 d_0 B_5)$.
\end{lemma}

\begin{proof}
There is a relation $P h_1 \cdot \D^3 h_1 h_3 = \tau \D^3 h_1^3 d_0$
in the Adams $E_2$-page.  Because of the differential
$d_2(\D^3 h_1 e_0) = \D^3 h_1^3 d_0 + \tau^5 e_0^2 g m$,
we have the relation
$P h_1 \cdot \D^3 h_1 h_3 = \tau^6 e_0^2 g m$ in
the $E_3$-page.

There is a differential $d_4(\tau^6 e_0^2 g m) = \tau^4 d_0^4 l$.
But $\tau^4 d_0^4 l$ is not divisible by $P h_1$,
so $\tau^6 e_0^2 g m$ cannot be divisible by $P h_1$
in the $E_4$-page.  Therefore,
$d_3(\D^3 h_1 h_3)$ must be non-zero.

The same argument shows that 
$d_3(\D^3 h_1 h_3 + \tau^3 d_0 B_5)$ must also be non-zero.
\end{proof}

\begin{remark}
In fact, Chua has determined that $d_3(\D^3 h_1 h_3)$ equals
$\tau^4 \D h_1 e_0^2 g$ \cite{Chua21}.
\end{remark}
}

\rev{
\begin{lemma}
\label{lem:d3-t^2d0B5}
\revdeg{80, 14, 42}
$d_3(\tau^2 d_0 B_5)$ equals either $\D^2 h_0 d_0 e_0$
or $\D^2 h_0 d_0 e_0 + \tau^3 \D h_1 e_0^2 g$.
\end{lemma}

\begin{proof}
The element $\D^2 h_0 d_0 e_0$ is a permanent cycle because
there are no possible differentials that it could support.
Moreover, 
it must map to zero under the inclusion of the bottom cell into $C\tau$ because there are no elements in the
Adams $E_\infty$-page for $C\tau$ of sufficiently high filtration.
Therefore, $\D^2 h_0 d_0 e_0$ 
is either hit by a differential, or it is the target of a 
hidden $\tau$ extension, or it is the target of a non-hidden
$\tau$ extension.

The only possible hidden $\tau$ extension has source $h_1^3 x_{76,6}$.
However, Table \ref{tab:hid-proj} shows that
$h_1^3 x_{76,6}$ is in the image of projection from $C\tau$ to the
top cell.  Therefore, it cannot support a hidden extension.

We now know that $\D^2 h_0 d_0 e_0$ must be hit by a differential, or
it is $\tau$-divisible in the $E_\infty$-page.
Lemma \ref{lem:perm-D^2d1} rules out one possible source for the
differential.  The only remaining possibilities are that 
$d_3(\tau^2 d_0 B_5)$ equals
$\D^2 h_0 d_0 e_0$ or 
$\D^2 h_0 d_0 e_0 + \tau^3 \D h_1 e_0^2 g$.
\end{proof} 

\begin{remark}
We are grateful to Dexter Chua for pointing out an error in a previous
version of Lemma \ref{lem:d3-t^2d0B5}.  In fact, Chua has determined
that the value of $d_3(\tau^2 d_0 B_5)$ is 
$\D^2 h_0 d_0 e_0 + \tau^3 \D h_1 e_0^2 g$ \cite{Chua21}.
\end{remark}
}

\begin{lemma}
\label{lem:d3-h2h4h6}
\mbox{}
\begin{enumerate}
\item
\revdeg{81, 3, 42}
$d_3(h_2 h_4 h_6) = 0$.
\item
\revdeg{82, 10, 42}
$d_3(P^2 h_2 h_6) = 0$.
\end{enumerate}
\end{lemma}

\begin{proof}
The value of $d_3(h_2 h_4 h_6)$ is not $h_2 h_6 d_0$ nor
$h_2 h_6 d_0 + \tau h_1 x_1$ by comparison to the Adams
spectral sequence for $C\tau$.

It remains to show that $d_3(h_2 h_4 h_6)$ cannot equal $\tau h_1 x_1$.
Suppose that the differential did occur.
Then there would be no possible targets for a 
hidden $\tau$ extension on $h_1 x_1$,
so the $\eta$ extension from $h_1 x_1$ to $h_1^2 x_1$ would be detected
by projection from $C\tau$ to the top cell.  But there 
is no such $\eta$ extension in the homotopy groups of $C\tau$.
This establishes the first formula.

The proof of the second formula is essentially the same,
using that the $\eta$ extension
from $\D^2 h_1 d_1$ to $\D^2 h_1^2 d_1$ cannot be
detected by projection from $C\tau$ to the top cell.
\end{proof}

\begin{lemma}
\label{lem:d3-D^2p}
\revdeg{81, 12, 42}
$d_3(\D^2 p) = 0$.
\end{lemma}

\begin{proof}
Suppose that $d_3(\D^2 p)$ were equal to $\tau^3 h_1 M e_0^2$.
In the Adams $E_4$-page,
the Massey product $\langle \tau^2 M g, \tau h_1 d_0, d_0 \rangle$
would equal $\tau^2 M \D h_2^2 g$, with no indeterminacy,
because of the Adams differential $d_3(\D h_2^2) = \tau h_1 d_0^2$
and because $d_0 \cdot \D^2 p = 0$.
By Moss's higher Leibniz rule \ref{thm:Leibniz},
$d_4(\tau^2 M \D h_2^2 g)$ would be a linear
combination of multiples of $\tau^2 M g$ and $d_0$.
But Table \ref{tab:Adams-d4} shows that $d_4(\tau^2 M \D h_2^2 g)$
equals $M P \D h_0^2 e_0$, which is not such a linear combination
in the Adams $E_4$-page.
\end{proof}

\begin{lemma}
\label{lem:d3-th6g}
\revdeg{83, 5, 43}
$d_3(\tau h_6 g + \tau h_2 e_2) = 0$.
\end{lemma}

\begin{proof}
In the Adams $E_3$-page, we have the matric Massey product
\[
\tau h_6 g + \tau h_2 e_2 =
\left\langle 
\left[
\begin{array}{cc}
\tau g & \tau h_2
\end{array}
\right],
\left[
\begin{array}{c}
h_5^2 \\ x_1
\end{array}
\right],
h_0
\right\rangle
\]
because of the Adams differentials
$d_2(h_6) = h_0 h_5^2$ and
$d_2(e_2) = h_0 x_1$, as well as the relation
$\tau g \cdot h_5^2 + \tau h_2 x_1$ in the Adams $E_2$-page.
Moss's higher Leibniz rule \ref{thm:Leibniz}
implies that 
$d_3(\tau h_6 g + \tau h_2 e_2)$ 
belongs to
\[
\left\langle
\left[
0 \enskip 0
\right],
\left[
\begin{array}{c}
h_5^2 \\ x_1
\end{array}
\right],
h_0 \right\rangle +
\left\langle
\left[
\tau g \enskip \tau h_2
\right],
\left[
\begin{array}{c}
0 \\ \tau h_1 m_1
\end{array}
\right],
h_0 \right\rangle +
\left\langle
\left[
\tau g \enskip \tau h_2
\right],
\left[
\begin{array}{c}
h_5^2 \\ x_1
\end{array}
\right],
0 \right\rangle
\]
since $d_3(x_1) = \tau h_1 m_1$,
where the Massey products are formed in the Adams $E_3$-page using
the $d_2$ differential.
This expression simplifies to
$\left\langle
\left[
\tau g \enskip \tau h_2
\right],
\left[
\begin{array}{c}
0 \\ \tau h_1 m_1
\end{array}
\right],
h_0 \right\rangle$,
which equals $\{0, \tau h_0^2 h_4 Q_3 \}$.

Table \ref{tab:nu-extn} shows that there is a hidden $\nu$
extension from $h_0^2 h_4 Q_3$ to $P h_1 x_{76,6}$.
The element $\tau P h_1 x_{76,6}$ is non-zero in the Adams
$E_\infty$-page.  Therefore,
$h_0^2 h_4 Q_3$ supports a (hidden or not hidden) 
$\tau$ extension whose target is in Adams filtration at most $10$.
The only possibility is that
$\tau h_0^2 h_4 Q_3$ is non-zero in the Adams 
$E_\infty$-page.
\end{proof}

\rev{
\begin{lemma}
\label{lem:d3-f2}
\revdeg{84, 4, 44}
$d_3(f_2) = \tau h_1 h_4 Q_3$.
\end{lemma}

\begin{proof}
Table \ref{tab:nu-extn} shows that 
$\tau h_1 Q_3$ detects $\nu^2 \theta_5$, and 
Table \ref{tab:Toda} shows that
$\tau h_1 h_4 Q_3$ detects
$\langle \nu^2 \theta_5, 2, \sigma^2 \rangle$,
with indeterminacy in strictly higher Adams filtration.
This bracket contains
$\nu^2 \langle \theta_5, 2, \sigma \rangle$, so
$\tau h_1 h_4 Q_3$ detects a multiple of $\nu^2$.
The only possibility is that $\tau h_1 h_4 Q_3$ is a multiple
of $h_2^2$ in the $E_\infty$-page.
This implies that 
$d_3(f_2)$ equals $\tau h_1 h_4 Q_3$ or $\tau h_1 h_4 Q_3 + h_0^2 h_6 g$.
The latter possibility is ruled out by comparison to
$C\tau$.
\end{proof}
}

\rev{
\begin{lemma}
\label{lem:d3-tx85,6+h0^3c3}
\revdeg{85, 6, 45}
$d_3(\tau x_{85,6} + h_0^3 c_3) = 0$.
\end{lemma}

\begin{proof}
Let $\alpha$ be an element of $\pi_{66,35}$ that is detected by
$\tau h_2 C'$.  Then $\nu \alpha$ is detected by $\tau h_2^2 C'$,
and $\tau \nu \alpha$ is zero.

Let $\ol{\alpha}$
be an element of $\pi_{70,36} C\tau$ that is detected by $h_0^2 h_3 h_6$.
Projection from $C\tau$ to the top cell takes $\ol{\alpha}$ to
$\nu \alpha$.
Moreover, in the homotopy of $C\tau$, the Toda bracket
$\langle 2, \sigma^2, \ol{\alpha} \rangle$
is detected by $h_0^3 c_3$.

Now projection from $C\tau$ to the top cell takes
$\langle 2, \sigma^2, \ol{\alpha} \rangle$ to 
$\langle 2, \sigma^2, \nu \alpha \rangle$, which equals zero
by Lemma \ref{lem:2,sigma^2,th2^2C'}.
Therefore, $h_0^3 c_3$ maps to zero under projection to the top cell
of $C\tau$,
so it must be in the image of inclusion of the bottom cell.

There are two possibilities.
First, $\tau x_{85,6} + h_0^3 c_3$ could survive, and it could
map to $h_0^3 c_3$ under inclusion of the bottom cell of $C\tau$.
Second,
$\tau h_1 f_2$ could map to $h_0^3 c_3$ under inclusion of the 
bottom cell.  This could only occur if $d_{10}(h_1 f_2)$ equaled
$M \D h_1 d_0$
and $d_9(\tau x_{85,6} + h_0^3 c_3)$ equaled $\tau M \D h_1 d_0$.

In either case, $d_3(\tau x_{85,6} + h_0^3 c_3)$ is zero.
\end{proof}
}

\rev{
\begin{remark}
\label{rem:d3-tx85,6+h0^3c3}
In the proof of Lemma \ref{lem:d3-tx85,6+h0^3c3}, we have used that
$d_5(\tau p_1 + h_0^2 h_3 h_6)$ equals $\tau^2 h_2^2 C'$ in order
to conclude that $\tau \nu \alpha$ is zero.  This differential
depends on work in preparation \cite{BIX}.  

However, we can
also prove Lemma \ref{lem:d3-tx85,6+h0^3c3} independently of \cite{BIX}.
Lemma \ref{lem:d5-tp1+h0^2h3h6} shows that the other possible value
of $d_5(\tau p_1 + h_0^2 h_3 h_6)$ is 
$\tau^2 h_2^2 C' + \tau h_3(\D e_1 + C_0)$.  In this case,
let $\beta$ be an element of $\pi_{62,33}$ that is detected by
$\D e_1 + C_0$.  Then $\nu \alpha + \sigma \beta$ is detected by 
$\tau h_2^2 C' + h_3(\D e_1 + C_0)$, 
and $\tau (\nu \alpha + \sigma \beta)$ is zero.

Projection from $C\tau$ to the top cell takes 
$\ol{\alpha}$ to $\nu \alpha + \sigma \beta$, and takes
$\langle 2, \sigma^2, \ol{\alpha} \rangle$ to 
$\langle 2, \sigma^2, \nu \alpha + \sigma \beta \rangle$, 
which equals zero
by Lemmas \ref{lem:2,sigma^2,th2^2C'} and 
\ref{lem:2,sigma^2,h3(De1+C0)}.
As in the proof of Lemma \ref{lem:d3-tx85,6+h0^3c3},
$h_0^3 c_3$ maps to zero under projection to the top cell of $C\tau$,
so it must be in the image of inclusion of the bottom cell.
\end{remark}
}

\rev{
\begin{lemma}
\label{lem:d3-t^2Mh0d0k}
\revdeg{88, 18, 46}
$d_3(\tau^2 M h_0 d_0 k) = P \D^2 h_0 d_0 e_0 + \tau^3 \D h_1 d_0^2 e_0^2$.
\end{lemma}

\begin{proof}
Table \ref{tab:Adams-d4} shows that
$d_4(\tau^2 M h_1 e_0) = M P^2 h_1$.
Multiply by $\tau d_0^2$ to see that 
$d_4(\tau^3 M h_1 d_0^2 e_0) = \tau M P^2 h_1 d_0^2$.
We have the relation $h_2 \cdot \tau^2 M h_0 d_0 k = \tau^3 M h_1 d_0^2 e_0$,
but $\tau M P^2 h_1 d_0^2$ is not divisible by $h_2$.
Therefore,
$\tau^2 M h_0 d_0 k$ cannot survive to the $E_4$-page.
By comparison to $\mmf$, there is only one possible value for
$d_3(\tau^2 M h_0 d_0 k)$.
\end{proof}

\begin{remark}
We are grateful to Dexter Chua for pointing out a small error in a previous
version of Lemma \ref{lem:d3-t^2Mh0d0k}.
\end{remark}
}

\begin{lemma}
\label{lem:d3-h2B5g}
\revdeg{89, 15, 50}
$d_3(h_2 B_5 g) = M h_1 c_0 e_0^2$.
\end{lemma}

\begin{proof}
\revv{We will first establish the relation
$d_0 \cdot h_2 B_5 g = \tau M h_1 e_0 g^2$.
We use the map $p_*$ of \cite{Isaksen19}*{Theorem 1.1}.
We have that $p_*(d_0 \cdot h_2 B_5 g) = p_*(\tau M h_1 e_0 g^2)$.
Therefore, 
$d_0 \cdot h_2 B_5 g$ equals $\tau M h_1 e_0 g^2$, modulo a possible
error term $P h_1^7 h_6 c_0 e_0$ in the kernel of $p_*$.
However, multiplication by $h_1$ eliminates the error term.
}

Table \ref{tab:Adams-d3} shows that $d_3(\tau e_0 g^2) = c_0 d_0 e_0^2$.
Therefore, $d_3(\tau M h_1 e_0 g^2)$ equals
$M h_1 c_0 d_0 e_0^2$.  Observing
that $M h_1 c_0 d_0 e_0^2$ is in fact non-zero in the Adams $E_3$-page,
we conclude that $d_3(h_2 B_5 g)$ must equal $M h_1 c_0 e_0^2$.
\end{proof}

\begin{lemma}
\label{lem:perm-M^2}
\revdeg{90, 8, 48}
The element $M^2$ is a permanent cycle.
\end{lemma}

\begin{proof}
Table \ref{tab:Massey} shows that the Massey product 
$\langle M h_1, h_0, h_0^2 g_2 \rangle$ equals $M^2 h_1$.
Therefore,
$M^2 h_1$ detects the Toda bracket
$\langle \eta \theta_{4.5}, 2, \sigma^2 \theta_4 \rangle$.
The indeterminacy consists entirely
of multiples of $\eta \theta_{4.5}$.
The Toda bracket contains
$\theta_4 \langle \eta \theta_{4.5}, 2, \sigma^2 \rangle$.
Now 
$\langle \eta \theta_{4.5}, 2, \sigma^2 \rangle$ is zero because
$\pi_{61,33}$ is zero.

We have now shown that $M^2 h_1$ detects a multiple
of $\eta$.  In fact, it detects a non-zero multiple of $\eta$
because $M^2 h_1$ cannot be hit by a differential by 
comparison to the Adams spectral sequence for $C\tau$.

Therefore, there exists a non-zero element of $\pi_{90,48}$
that is detected in Adams fitration at most $12$.
The only possibility is that $M^2$ survives.
\end{proof}

\begin{lemma}
\label{lem:d3-Dh2^2h6}
\revdeg{93, 7, 48}
$d_3(\D h_2^2 h_6) = \tau h_1 h_6 d_0^2$.
\end{lemma}

\begin{proof}
In the Adams $E_3$-page, 
$\D h_2^2 h_6$ equals $\langle \D h_2^2, h_5^2, h_0 \rangle$,
with no indeterminacy, because of the Adams differential
$d_2(h_6) = h_0 h_5^2$.
Using that $d_3(\D h_2^2) = \tau h_1 d_0^2$, 
Moss's higher Leibniz rule \ref{thm:Leibniz}
implies that $d_3(\D h_2^2 h_6)$ is contained in
\[
\langle \tau h_1 d_0^2, h_5^2, h_0 \rangle +
\langle \D h_2^2, 0, h_0 \rangle +
\langle \D h_2^2, h_5^2, 0 \rangle.
\]
All of these brackets have no indeterminacy, and the last two
equal zero.
The first bracket equals $\tau h_1 h_6 d_0^2$, using the
Adams differential $d_2(h_6) = h_0 h_5^2$.
\end{proof}

\begin{lemma}
\label{lem:d3-P^2h6d0}
\revdeg{93, 13, 48}
$d_3(P^2 h_6 d_0) = 0$.
\end{lemma}

\begin{proof}
In the Adams $E_3$-page, the element $P^2 h_6 d_0$ equals
the Massey product
$\langle P^2 d_0, h_5^2, h_0 \rangle$, with no indeterminacy,
because of the Adams differential $d_2(h_6) = h_0 h_5^2$.
Moss's higher Leibniz rule \ref{thm:Leibniz} implies
that $d_3(P^2 h_6 d_0)$ is a linear combination of multiples of
$h_0$ and of $P^2 d_0$.  The only possibility is that
$d_3(P^2 h_6 d_0)$ is zero.
\end{proof}

\rev{
\begin{lemma}
\label{lem:d3-t^2MPh0d0j}
\revdeg{93, 22, 48}
$d_3(\tau^2 M P h_0 d_0 j) = P^2 \D^2 h_0 d_0^2 + 
\tau^3 P \D h_1 d_0^3 e_0$.
\end{lemma}
}
\begin{proof}
Table \ref{tab:Adams-d4} shows that
$d_4(\tau^2 P d_0 e_0) = P^3 d_0$.
Multiplication by $\tau M P h_1$ shows that
$d_4(\tau^3 M P^2 h_1 d_0 e_0)$ equals
$\tau M P^4 h_1 d_0$.
But $\tau^3 M P^2 h_1 d_0 e_0$ equals
$h_0 \cdot \tau^2 M P h_0 d_0 j$, 
while $\tau M P^4 h_1 d_0$ is not divisible by $h_0$.
Therefore,
$\tau^2 M P h_0 d_0 j$ cannot survive to the $E_4$-page.

The possible values for
$d_3(\tau^2 M P h_0 d_0 j)$ are
\revv{the non-zero}
linear combinations of 
$P^2 \D^2 h_0 d_0^2$ and
$\tau^3 P \D h_1 d_0^3 e_0$.
\revv{The map to the Adams spectral sequence for $\tmf$
takes both $\tau^2 M P h_0 d_0 j$ and 
$P^2 \D^2 h_0 d_0^2 + \tau^3 P \D h_1 d_0^3 e_0$ to zero, 
but it takes each of
$P^2 \D^2 h_0 d_0^2$ and
$\tau^3 P \D h_1 d_0^3 e_0$ to the unique non-zero element
in the appropriate degree.
Therefore,
$d_3(\tau^2 M P h_0 d_0 j)$ cannot equal either
$P^2 \D^2 h_0 d_0^2$ or $\tau^3 P \D h_1 d_0^3 e_0$.
}
\end{proof}

\begin{lemma}
\label{lem:perm-P^3h6c0}
\revdeg{95, 16, 49}
The element $P^3 h_6 c_0$ is a permanent cycle.
\end{lemma}

\begin{proof}
Table \ref{tab:eta-extn} shows that
$P^3 c_0$ detects the product $\eta \rho_{31}$.
Using the Moss Convergence Theorem \ref{thm:Moss} 
and the Adams differential $d_2(h_6) = h_0 h_5^2$,
the element $P^3 h_6 c_0$ must survive to detect the
Toda bracket
$\langle \eta \rho_{31}, 2, \theta_5 \rangle$.
\end{proof}

\begin{remark}
\label{rem:perm-P^3h6c0}
We suspect that $P^3 h_6 c_0$ detects the product
$\eta_6 \rho_{31}$.  However, the argument of
Lemma \ref{lem:rho15-h1h6} cannot be completed because
the Toda bracket
$\langle \eta \rho_{31}, 2, \theta_5 \rangle$ might
have indeterminacy in lower Adams filtration.
\end{remark}

\section{The Adams $d_4$ differential}
\label{sctn:Adams-d4}

\revv{
Table \ref{tab:Adams-d4} lists the multiplicative generators 
of the Adams $E_4$-page through the 95-stem
whose $d_4$ differentials are non-zero, or whose $d_4$ differentials
are zero for non-obvious reasons.
}

\begin{thm}
\label{thm:Adams-d4}
Table \ref{tab:Adams-d4} lists some values of the
Adams $d_4$ differential on multiplicative generators.
Through the 95-stem, 
the Adams $d_4$ differential is zero on all multiplicative
generators not listed in the table.
\end{thm}

\begin{proof}
The $d_4$ differential on many multiplicative generators is zero.
A few of these multiplicative generators appear
in Table \ref{tab:Adams-d4} because their proofs require further
explanation.
For the remaining majority of such multiplicative generators,
the $d_4$ differential is zero
because there are no possible non-zero values, 
or because of comparison to the Adams spectral sequences for $C\tau$,
$\tmf$, or $\mmf$.
In a few cases, the multiplicative generator
is already known to be a permanent cycle as shown in
Table \ref{tab:Adams-perm}.
These cases do not appear in Table \ref{tab:Adams-d4}.

The last column of Table \ref{tab:Adams-d4} gives information on the
proof of each differential.
Most follow immediately by comparison to
the Adams spectral sequence for $C\tau$, or
by comparison to the classical Adams spectral sequence for $\tmf$, or
by comparison to the $\C$-motivic Adams spectral sequence for $\mmf$.

If an element is listed in the last column of Table \ref{tab:Adams-d4}, 
then the corresponding differential can be deduced from a straightforward
argument using a multiplicative relation.  For example, 
\[
d_4( d_0 \cdot \tau^2 e_0 g^2 ) =
d_4( e_0^2 \cdot \tau^2 e_0 g) = 
e_0^2 \cdot P d_0^2 = d_0^5,
\]
so $d_4(\tau^2 e_0 g^2)$ must equal $d_0^4$.

The remaining more difficult computations are carried out in the
following lemmas.
\end{proof}




\begin{lemma}
\label{lem:d4-th1.Dx}
\revdeg{62, 10, 32}
$d_4(\tau h_1 \cdot \D x) = \tau^2 \D h_2^2 d_0 e_0$.
\end{lemma}

\revv{This differential was previously proved in
\cite{WangXu17}*{Remark 11.2}.  We repeat the argument here
for completeness.
}

\begin{proof}
Table \ref{tab:Adams-d4} shows that $\tau^3 \D h_2^2 g^2$ supports a $d_4$ differential, and 
Table \ref{tab:Adams-perm} shows that 
$\tau \D^2 h_1^2 g + \tau^3 \D h_2^2 g^2$ is a permanent cycle.
Therefore,
$\tau \D^2 h_1^2 g$ also supports a $d_4$ differential.

On the other hand, we have
\[
h_1 \cdot \tau \D^2 h_1 g = P h_1 \cdot \D x =
\D x \langle h_1, h_0^3 h_3, h_0 \rangle.
\]
This expression equals $\langle h_1 \cdot \D x, h_0^3 h_3, h_0 \rangle$
by inspection of indeterminacies.
Therefore, the Toda bracket
$\langle \{ \tau h_1 \cdot \D x \}, 8 \sigma, 2 \rangle$
cannot be well-formed, since otherwise it would be detected
by $\tau \D^2 h_1^2 g$.
The only possibility is that 
$\tau h_1 \cdot \D x$ is not a permanent cycle, and the
only possible differential is that
$d_4(\tau h_1 \cdot \D x)$ equals $\tau^2 \D h_2^2 d_0 e_0$.
\end{proof}

\begin{lemma}
\label{lem:d4-D^2h3^2}
\revdeg{62, 10, 32}
$d_4(\D^2 h_3^2) = 0$.
\end{lemma}

\begin{proof}
Table \ref{tab:Adams-d5} shows that 
$d_5(\tau h_1^2 \cdot \D x)$ equals
$\tau^3 d_0^2 e_0^2$.  
The element $\tau^3 d_0^2 e_0^2$ is not divisible by
$h_1$ in the $E_5$-page, so
$\tau h_1^2 \cdot \D x$ 
cannot be divisible by $h_1$ in the $E_4$-page.

If $d_4(\D^2 h_3^2)$ equaled 
$\tau^2 \D h_2^2 d_0 e_0$, then
$\D^2 h_3^2 + \tau h_1 \cdot \D x$ would survive to the $E_5$-page, 
and $\tau^2 h_1^2 \cdot \D x$ would be divisible by $h_1$ in the $E_5$-page.
\end{proof}

\begin{lemma}
\label{lem:d4-tX2}
\revdeg{63, 7, 33}
$d_4(\tau X_2) = \tau M h_2 d_0$.
\end{lemma}

\begin{proof}
Table \ref{tab:Adams-d4} shows that $d_4(C')$ equals $M h_2 d_0$.
Therefore, either $\tau X_2$ or $\tau X_2 + \tau C'$ is non-zero on the
$E_\infty$-page.  The inclusion of the bottom cell into $C\tau$
takes this element to $\ol{h_5 d_0 e_0}$.

In the homotopy of $C\tau$, there is a $\nu$ extension from
$\ol{h_5 d_0 e_0}$ to $\tau B_5$, and
inclusion of the bottom cell into $C\tau$ takes
$\tau h_2 C'$ to $\tau B_5$.

It follows that there must be a $\nu$ extension with target
$\tau h_2 C'$.  The only possibility is that
$\tau X_2 + \tau C'$ is non-zero on the $E_\infty$-page,
and therefore $d_4(\tau X_2)$ equals $d_4(\tau C')$.
\end{proof}

\begin{lemma}
\label{lem:d4-h0d2}
\revdeg{68, 5, 36}
$d_4(h_0 d_2) = X_3$.
\end{lemma}

\begin{proof}
The element $X_3$ is a permanent cycle.  The only possible target for 
a differential is $\tau^2 d_0 e_0 m$, but this is ruled out by 
comparison to $\tmf$.

\revv{
In the Adams $E_\infty$-page for $C\tau$, there are several elements in stem 67 and weight 36.  However, they all have filtration less than 9.
Since $X_3$ has filtration $9$, it must map to zero under inclusion
of the bottom cell into $C\tau$.}
Therefore, $X_3$ is the target of a hidden
$\tau$ extension, or it is hit by a differential.

The only possible hidden $\tau$ extension would have source
$h_1 \cdot \D_1 h_3^2$.  In $C\tau$, there is an $\eta$ extension from
$h_0 d_2$ to $\ol{h_1^2 \cdot \D_1 h_3^2}$.  
Since $\ol{h_1^2 \cdot \D_1 h_3^2}$ maps non-trivially 
(to $h_1^2 \cdot \D_1 h_3^2$) under projection to the top cell of
$C\tau$, it follows that
$h_0 d_2$ also maps non-trivially under projection.
\revv{For degree reasons,}
the only possibility is that $h_0 d_2$ maps to
$h_1 \cdot \D_1 h_3^2$, and therefore
$h_1 \cdot \D_1 h_3^2$ does not support a hidden 
$\tau$ extension.

Therefore, $X_3$ must be hit by a differential, and there is just
one possibility.
\end{proof}

\begin{lemma}
\label{lem:d4-Mh2g}
\revdeg{68, 11, 38}
$d_4(M h_2 g) = 0$.
\end{lemma}

\begin{proof}
Table \ref{tab:Massey} shows that the Massey product
$\langle h_2 g, h_0^3, g_2 \rangle$ equals $M h_2 g$.
The Moss Convergence Theorem \ref{thm:Moss} shows that
$M h_2 g$ must survive to detect the Toda bracket
$\langle \{h_2 g\}, 8, \kappabar_2 \rangle$.
\end{proof}

\begin{lemma}
\label{lem:d4-h2^2G0}
\revdeg{72, 9, 40}
$d_4(h_2^2 G_0) = \tau g^2 n$.
\end{lemma}

\begin{proof}
Table \ref{tab:2-extn} shows that
there is a hidden $2$ extension from 
$h_0 h_3 g_2$ to $\tau g n$.  Therefore, 
$\tau g n$ detects $4 \sigma \kappabar_2$.

Table \ref{tab:Massey} shows that
$\langle h_1^3 h_4, h_1, \tau g n \rangle$ consists of the
two elements
$\tau g^2 n$ and $\tau g^2 n + M h_2 g \cdot h_0 h_2$.
Then the Toda bracket
$\langle \eta^2 \eta_4, \eta, 4 \sigma \kappabar_2 \rangle$
is detected by either
$\tau g^2 n$ or $\tau g^2 n + M h_2 g \cdot h_0 h_2$.
But $M h_2 g \cdot h_0 h_2$ is hit by an Adams $d_2$ differential,
so $\tau g^2 n$ detects the Toda bracket.

The Toda bracket has no indeterminacy, so it equals
$\langle \eta^2 \eta_4, \eta, 2 \rangle 2 \sigma \kappabar_2$.
This last expression must be zero.

We have shown that $\tau g^2 n$ must be hit by some differential.
The only possibility is that $d_4(h_2^2 G_0) = \tau g^2 n$.
\end{proof}

\begin{lemma}
\label{lem:d4-Dh0^2h3g2}
\revdeg{75, 11, 40}
$d_4(\D h_0^2 h_3 g_2) = \tau M h_1 d_0^2$.
\end{lemma}

\begin{proof}
Table \ref{tab:Adams-d5} shows that 
$d_5(A') = \tau M h_1 d_0$.
Now $d_0 A'$ is zero in the $E_5$-page, so
$\tau M h_1 d_0^2$ must also be zero in the $E_5$-page.
\end{proof}

\begin{lemma}
\label{lem:d4-D^2h1h3g}
\revdeg{76, 14, 41}
$d_4(\D^2 h_1 h_3 g) = \tau \D h_2^2 d_0^2 e_0$.
\end{lemma}

\begin{proof}
Table \ref{tab:Toda} shows that the element
$\D h_2^2 d_0^2 e_0$ detects the Toda bracket
$\langle \tau \eta \kappa \kappabar^2, \eta, \eta^2 \eta_4 \rangle$.
Now shuffle to obtain
\[
\tau \langle \tau \eta \kappa \kappabar^2, \eta, \eta^2 \eta_4 \rangle =
\langle \tau, \tau \eta \kappa \kappabar^2, \eta \rangle \eta^2 \eta_4.
\]
Table \ref{tab:Toda} shows that
$\langle \tau, \tau \eta \kappa \kappabar^2, \eta \rangle$ is detected
by $h_0 h_2 h_5 i$.
It follows that the expression
$\langle \tau, \tau \eta \kappa \kappabar^2, \eta \rangle \eta^2 \eta_4$
is zero, so $\tau \D h_2^2 d_0^2 e_0$ must be hit by some differential.
The only possibility is that 
$d_4(\D^2 h_1 h_3 g)$ equals 
$\tau \D h_2^2 d_0^2 e_0$.
\end{proof}

\begin{lemma}
\label{lem:d4-h0e2}
\revdeg{80, 5, 42}
$d_4(h_0 e_2) = \tau h_1^3 x_{76,6}$.
\end{lemma}

\begin{proof}
Table \ref{tab:misc-extn}
shows that $\sigma^2 \theta_5$ is detected by
$h_0 h_4 A$ \rev{or $h_0 h_4 A + \tau^2 d_1 g_2$.
Note that both $h_2 \cdot h_0 h_4 A$ and 
$h_2 (h_0 h_4 A + \tau^2 d_1 g_2)$ equal 
$\tau h_1^3 x_{76,6}$.
}
Since $\nu \sigma = 0$, the element
$\tau h_1^3 x_{76,6}$ must be hit
by a differential.
The only possibility is that $d_4(h_0 e_2)$ equals
$\tau h_1^3 x_{76,6}$.
\end{proof}

\begin{lemma}
\label{lem:d4-D^3h1^2h3}
\revdeg{81, 15, 42}
$d_4(\D^3 h_1^2 h_3) = \tau^4 d_0 e_0^2 l$.
\end{lemma}

\begin{proof}
Table \ref{tab:eta-extn} shows that there is a hidden
$\eta$ extension from $\tau^2 \D h_1 g^2$ to $\tau^2 d_0 e_0 m$.
\revv{Multiply by $d_0$ to see that}
there is also a hidden $\eta$ extension
from $\tau^2 \D h_1 e_0^2 g$ to $\tau^2 d_0 e_0^2 l$.

Also, $\tau^2 \D h_1 e_0^2 g$ 
detects an element in $\pi_{79,43}$ that is annihilated by $\tau^2$.
Therefore, $\tau^4 d_0 e_0^2 l$ must be hit by some
differential.
Moreover, comparison to $\mmf$ shows that 
$\tau^3 d_0 e_0^2 l$ is not hit by a differential.

The hidden $\eta$ extension from $\tau^3 \D h_1 e_0^2 g$
to $\tau^3 d_0 e_0^2 l$ is detected by projection from
$C\tau$ to the top cell.
The only possibility is that this hidden $\eta$ extension is
the image of the $h_1$ extension
from $\D^3 h_1 h_3$ to $\D^3 h_1^2 h_3$
in the Adams $E_\infty$-page for $C\tau$.

Therefore, $\D^3 h_1^2 h_3$ maps non-trivially under projection
from $C\tau$ to the top cell.
Consequently,
$\D^3 h_1^2 h_3$ cannot be a permanent cycle in the
Adams spectral sequence for the sphere.
\end{proof}

\begin{lemma}
\label{lem:d4-Dj1}
\revdeg{83, 11, 45}
$d_4(\D j_1) = \tau M h_0 e_0 g$.
\end{lemma}

\begin{proof}
Otherwise, both $\D j_1$ and $\tau g C'$ would survive to the
$E_\infty$-page, and neither could be the target of a
hidden $\tau$ extension.  They would both map non-trivially
under inclusion of the bottom cell into $C\tau$.
But there are not enough elements in $\pi_{83,45} C\tau$
for this to occur.
\end{proof}

\rev{
\begin{lemma}
\label{lem:d4-h1f2}
\revdeg{85, 5, 45}
$d_4(h_1 f_2) = 0$.
\end{lemma}

\begin{proof}
Table \ref{tab:nu-extn} shows that there is a hidden
$\nu$ extension from $h_2 g D_3$ to $B_6 d_1$.
If $d_4(h_1 f_2)$ equaled $\tau h_2 g D_3$,
then this $\nu$ extension would be in the image of
projection from $C\tau$ to the top cell, since $h_2 g D_3$
cannot support a hidden $\tau$ extension.
However, there is no such $\nu$ extension in the homotopy 
of $C\tau$.
\end{proof}
}

\rev{
\begin{lemma}
\label{lem:d4-tx85,6+h0^3c3}
\revdeg{85, 6, 44}
$d_4(\tau x_{85,6} + h_0^3 c_3) = 0$.
\end{lemma}

\begin{proof}
We showed in Lemma \ref{lem:d3-tx85,6+h0^3c3} that
$h_0^3 c_3$ is in the image of inclusion of the bottom cell of $C\tau$.
Therefore, $P x_{76,6}$ cannot be in the image of projection from
$C\tau$ to the top cell.  Since $P x_{76,6}$ cannot support a hidden
$\tau$ extension, there can be no differential whose value is
$\tau P x_{76,6}$.
\end{proof}
}

\begin{lemma}
\label{lem:d4-h1c3}
\revdeg{86, 4, 45}
$d_4(h_1 c_3) = \tau h_0 h_2 h_4 Q_3$.
\end{lemma}

\begin{proof}
Lemma \ref{lem:compound-(epsilon+etasigma)-theta5}
shows that there exists an element
$\alpha$ in $\pi_{67,36}$ that is detected by
$h_0 Q_3 + h_0 n_1$ such that
$\tau \nu \alpha$ equals
$(\eta \sigma + \epsilon) \theta_5$.

Table \ref{tab:Toda} shows that the Toda bracket
$\langle \nu, \sigma, 2\sigma \rangle$ is detected by $h_2 h_4$, so
the element
$\tau h_0 h_2 h_4 Q_3$ detects
$\tau \alpha \langle \nu, \sigma, 2 \sigma \rangle$,
which is contained in
$\langle \tau \nu \alpha, \sigma, 2 \sigma \rangle$.
The indeterminacy in these expressions
is zero because
$\tau \nu \alpha \cdot \pi_{15,8}$ and
$2 \sigma \cdot \pi_{78,41}$ are both zero.

We now know that
$\tau h_0 h_2 h_4 Q_3$ detects the Toda bracket
$\langle (\epsilon + \eta \sigma) \theta_5, \sigma, 2 \sigma \rangle$.
This bracket contains
$\theta_5 \langle \epsilon + \eta \sigma, \sigma, 2 \sigma \rangle$.
Lemma \ref{lem:nubar,sigma,2sigma} shows that 
the bracket
$\langle \epsilon + \eta \sigma, \sigma, 2 \sigma \rangle$ contains
$0$, so
$\theta_5 \langle \epsilon + \eta \sigma, \sigma, 2 \sigma \rangle$
equals zero.

Finally, we have shown that $\tau h_0 h_2 h_4 Q_3$
detects zero, so it must be hit by some differential.
\end{proof}

\begin{lemma}
\label{lem:d4-x87,7}
\revdeg{87, 7, 45}
$d_4(x_{87,7}) = 0$.
\end{lemma}


\begin{proof}
Consider the exact sequence
\[
\pi_{87,45} \map \pi_{87,45} C\tau \map \pi_{86,46}.
\]
The middle term $\pi_{87,45} C\tau$ is isomorphic to $(\Z/2)^4$.
The elements of $\pi_{87,45}$ that are not divisible by
$\tau$ are detected by $P^2 h_6 c_0$, and possibly
$x_{87,7}$ and $\tau \D h_1 H_1$.
On the other hand, the elements of $\pi_{86,46}$ that are
annihilated by $\tau$ are detected by
$\tau^3 \D c_0 e_0^2 g$ and possibly $M \D h_0^2 e_0$.

In order for the possibility $M \D h_0^2 e_0$ to occur,
either $x_{87,7}$ or $\tau \D h_1 H_1$ would have to support
a differential hitting $\tau M \D h_0^2 e_0$, in which case
one of those possibilities could not occur.

If $d_4(x_{87,7})$ equaled $\tau^3 g G_0$, then there would not
be enough elements to make the above sequence exact.
\end{proof}

\begin{lemma}
\label{lem:d4-tDh1H1}
\revdeg{87, 10 45}
$d_4(\tau \D h_1 H_1) = 0$.
\end{lemma}

\begin{proof}
The element $\D^2 h_2^2 d_1$ is a permanent cycle
that cannot be hit by any differential because
$h_2 \cdot \D^2 h_2^2 d_1$ cannot be hit by a differential.
The element $\D^2 h_2^2 d_1$ cannot be in the image
of projection from $C\tau$ to the top cell,
and it cannot support a hidden $\tau$ extension.
Therefore,
$\tau \D^2 h_2^2 d_1$ cannot be hit by a differential.
\end{proof}

\begin{lemma}
\label{lem:d4-th2B5g}
\revdeg{89, 15, 49}
$d_4(\tau h_2 B_5 g) = M h_1 d_0^3$.
\end{lemma}

\begin{proof}
Table \ref{tab:misc-extn} shows that
$M d_0$ detects $\kappa \theta_{4.5}$.
Therefore, $M d_0^3$ detects $\kappa^3 \theta_{4.5}$,
which equals $\eta^2 \kappabar^2 \theta_{4.5}$ because
Table \ref{tab:eta-extn} shows that there is a hidden
$\eta$ extension from $\tau^2 h_1 g^2$ to $d_0^3$.

Now $\eta^2 \kappabar^2 \theta_{4.5}$ is zero because
$\eta^2 \kappabar \theta_{4.5}$ is zero.
Therefore, $M d_0^3$ and $M h_1 d_0^3$ must both be hit
by differentials.

There are several possible differentials that can hit
$M h_1 d_0^3$.  
The element $h_1 x_{88,10}$ cannot be the source of this differential
because Table \ref{tab:Adams-perm} shows that
$x_{88,10}$ is a permanent cycle.
The element $\tau h_2^2 g C'$ cannot be the source of the differential
because $h_2^2 g C'$ is a permanent cycle by comparison to $\mmf$.
The element $\D h_1 g_2 g$ cannot be the source 
because it equals $h_3 (\D e_1 + C_0) g$.
The only remaining possibility is that
$d_4(\tau h_2 B_5 g)$ equals $M h_1 d_0^3$.
\end{proof}

\begin{lemma}
\label{lem:d4-Dh2^2A'}
\revdeg{91, 12, 48}
$d_4(\D h_2^2 A') = 0$.
\end{lemma}

\begin{proof}
In the Adams $E_4$-page, the element $\D h_2^2 A'$ equals the 
Massey product $\langle A', h_1, \tau d_0^2 \rangle$,
with no indeterminacy
because of the Adams differential $d_3(\D h_2^2) = \tau h_1 d_0^2$.
Moss's higher Leibniz rule \ref{thm:Leibniz}
implies that $d_4(\D h_2^2 A')$ is contained in
\[
\langle 0, h_1, \tau d_0^2 \rangle + \langle A', 0, \tau d_0^2 \rangle +
\langle A', h_1, 0 \rangle,
\]
so it is
a linear combination
of multiples of $A'$ and $\tau d_0^2$.
The only possibility is that $d_4(\D h_2^2 A')$ is zero.
\end{proof}

\begin{lemma}
\label{lem:d4-h4^2h6}
\revdeg{93, 3, 48}
$d_4(h_4^2 h_6) = h_0^3 g_3$.
\end{lemma}

\begin{proof}
By comparison to 
the Adams spectral sequence for $C\tau$, the value of 
$d_4(h_4^2 h_6)$ is either
$h_0^3 g_3$ or $h_0^3 g_3 + \tau h_1 h_4^2 D_3$.

Table \ref{tab:misc-extn} shows that 
$h_0^2 g_3$ detects the product $\theta_4 \theta_5$.
Since $2 \theta_4 \theta_5$ equals zero, 
$h_0^3 g_3$ must be hit by a differential.
\end{proof}

\rev{
\begin{lemma}
\label{lem:d4-MD^2h1^2}
\revdeg{95, 16, 50}
$d_4(M \D^2 h_1^2) = M P \D h_0^2 e_0$.
\end{lemma}

\begin{proof}
Table \ref{tab:Toda} shows that
$M \D^2 h_1^2 + \tau^2 M \D h_2^2 g$ detects the Toda bracket
$\langle \eta, \tau \kappa^2, \tau \theta_{4.5} \kappabar \rangle$.
Therefore,
$d_4(M \D^2 h_1^2)$ equals $d_4(\tau^2 M \D h_2^2 g)$.
\end{proof}
}

\section{The Adams $d_5$ differential}
\label{sctn:Adams-d5}

\revv{
Table \ref{tab:Adams-d5} lists the multiplicative generators 
of the Adams $E_5$-page through the \rev{92}-stem
whose $d_5$ differentials are non-zero, or whose $d_5$ differentials
are zero for non-obvious reasons.
}

\begin{thm}
\label{thm:Adams-d5}
Table \ref{tab:Adams-d5} lists some values of the
Adams $d_5$ differential on multiplicative generators.
Through the \rev{92}-stem, 
the Adams $d_5$ differential is zero on all multiplicative
generators not listed in the table.
\end{thm}

\begin{proof}
The $d_5$ differential on many multiplicative generators is zero.
For the majority of such multiplicative generators,
the $d_5$ differential is zero
because there are no possible non-zero
values, or by comparison to the Adams spectral sequence for $C\tau$,
or by comparison to $\tmf$ or $\mmf$.
In a few cases, the multiplicative generator
is already known to be a permanent cycle; $h_1 h_6$ is one
such example.
A few additional cases appear
in Table \ref{tab:Adams-d5} because their proofs require further
explanation.

The last column of Table \ref{tab:Adams-d5} gives information on the
proof of each differential.
Many computations follow immediately by comparison to
the Adams spectral sequence for $C\tau$.

If an element is listed in the last column of Table \ref{tab:Adams-d5}, 
then the corresponding differential can be deduced from a straightforward
argument using a multiplicative relation.  For example, 
\[
d_5( \tau \cdot g A' ) =
d_5( \tau g \cdot A') =
\tau g \cdot \tau M h_1 d_0 = \tau^2 M h_1 e_0^2,
\]
so $d_5(g A')$ must equal $\tau M h_1 e_0^2$.

A few of the more difficult computations appear in \cite{BIX}.
The remaining more difficult computations are carried out in the
following lemmas.
\end{proof}




\begin{lemma}
\label{lem:d5-th1^2.Dx}
\revdeg{63, 11, 33}
$d_5(\tau h_1^2 \cdot \D x) = \tau^3 d_0^2 e_0^2$.
\end{lemma}

\begin{proof}
The element $\tau^2 d_0^2 e_0^2$ cannot be hit by a differential.
There is a hidden $\eta$ extension
from $\tau \D h_2^2 d_0 e_0$ to $\tau^2 d_0^2 e_0^2$
because of the hidden $\tau$ extensions
from $\tau h_1 g^3 + h_1^5 h_5 c_0 e_0$ to
$\D h_2^2 d_0 e_0$ and
from $h_1^6 h_5 c_0 e_0$ to $d_0^2 e_0^2$.
This shows that $\tau^3 d_0^2 e_0^2$ must be hit by some differential.

This hidden $\eta$ extension is detected by
projection from $C\tau$ to the top cell.
Since $\ol{P h_5 c_0 d_0}$ in $C\tau$ maps to
$\tau \D h_2^2 d_0 e_0$ under projection to the top cell,
it follows that $\ol{P h_1 h_5 c_0 d_0}$ in $C\tau$
maps to $\tau^2 d_0^2 e_0^2$ under projection to the top cell.

If $\tau h_1^2 \cdot \D x$ survived, then it could not be the target of
a hidden $\tau$ extension and it could not be hit by a differential.
Also, it could not map non-trivially under inclusion of
the bottom cell into $C\tau$, since the only possible value
$\ol{P h_1 h_5 c_0 d_0}$ has already been accounted for 
in the previous paragraph.
\end{proof}

\begin{lemma}
\label{lem:d5-h5d0i}
\revdeg{68, 12, 36}
$d_5(h_5 d_0 i) = \tau \D h_1 d_0^3$.
\end{lemma}

\begin{proof}
We showed in Lemma \ref{lem:d3-tDg2.h0^3}
that $P h_2 h_5 j$ cannot be divisible by $h_0$ in the $E_\infty$-page.
Therefore, $h_5 d_0 i$ must support a differential.
\end{proof}

\begin{lemma}
\label{lem:d5-tp1+h0^2h3h6}
\revdeg{70, 4, 36}
$d_5(\tau p_1 + h_0^2 h_3 h_6)$ equals either
$\tau^2 h_2^2 C'$ or $\tau^2 h_2^2 C' + \tau h_3(\D e_1 + C_0)$.
\end{lemma}

\begin{proof}
Projection to the top cell of $C\tau$ takes
$h_4 D_2$ to $\tau^3 d_1 g^2$.
Moreover,
there is a $\nu$ extension in the homotopy of $C\tau$
from $h_0^2 h_3 h_6$ to $h_4 D_2$.
Therefore, this $\nu$ extension must be in the image of
projection to the top cell.

Table \ref{tab:nu-extn} shows that there is a hidden
$\nu$ extension from
$\tau h_2^2 C'$ to $\tau^3 d_1 g^2$.
Therefore, either $\tau h_2^2 C'$ or
$\tau h_2^2 C' + h_3(\D e_1 + C_0)$ is in the image of
projection to the top cell, so
$\tau^2 h_2^2 C'$ or $\tau^2 h_2^2 C' + \tau h_2 (\D e_1 + C_0)$
is hit by a differential.
The element $\tau p_1 + h_0^2 h_3 h_6$ is the 
only possible source for this differential.
\end{proof}

\begin{lemma}
\label{lem:d5-h1x71,6}
\revdeg{72, 7, 39}
$d_5(h_1 x_{71,6}) = 0$.
\end{lemma}

\begin{proof}
Table \ref{tab:tau-extn} shows that 
there is a hidden $\tau$ extension from 
$M h_1^2 h_3 g$ to $M h_1 d_0^2$.
Therefore, $M h_2^2 g$ must also support a $\tau$ extension.
This shows that $\tau M h_2^2 g$ cannot be the target of a differential.
\end{proof}

\begin{lemma}
\label{lem:d5-h4D2}
\revdeg{73, 7, 38}
$d_5(h_4 D_2) = \tau^4 d_1 g^2$.
\end{lemma}

\begin{proof}
Suppose for sake of contradiction that $h_4 D_2$ survived, and
let $\alpha$ be an element of $\pi_{73,38}$ that is detected by it.
Table \ref{tab:tau-extn} shows that there is a hidden $\tau$ extension
from $h_1^2 h_6 c_0$ to $h_0 h_4 D_2$.  Therefore,
$h_0 h_4 D_2$ detects both $2 \alpha$ and $\tau \eta \epsilon \eta_6$.
However, it is possible that the difference between these two elements
is detected by $\tau^2 M d_0^2$ or by $\tau^3 \D h_1 d_0 e_0^2$.
We will handle of each of these cases.

First, suppose that $2 \alpha$ equals $\tau \eta \epsilon \eta_6$.
Then the Toda bracket
\[
\left\langle
\eta,
\left[
\begin{array}{cc}
2 & \tau \eta \epsilon 
\end{array}
\right],
\left[
\begin{array}{c}
\alpha \\ \eta_6 
\end{array}
\right]
\right\rangle
\]
is well-defined.  Inclusion of the bottom cell into $C\tau$ takes this
bracket to
\[
\left\langle
\eta,
\left[
\begin{array}{cc}
2 & 0
\end{array}
\right],
\left[
\begin{array}{c}
\alpha \\ \eta_6 
\end{array}
\right]
\right\rangle =
\langle \eta, 2, \alpha \rangle,
\]
so $\langle \eta, 2, \alpha \rangle$ is in the image of inclusion of 
the bottom cell.

On the other hand,
in the homotopy of $C\tau$, the bracket 
$\langle \eta, 2, \alpha \rangle$ is detected by
$\ol{h_1^3 h_6 c_0}$, with indeterminacy generated by
$\ol{h_1^2 h_4 Q_2}$.  
These elements map non-trivially under projection to the top cell,
which contradicts that they are in the image of inclusion of the 
bottom cell.  

Next, suppose that $2 \alpha + \tau \eta \epsilon \eta_6$
is detected by $\tau^3 \D h_1 d_0 e_0^2$.
Then the Toda bracket
\[
\left\langle
\eta,
\left[
\begin{array}{ccc}
2 & \tau \eta \epsilon & \tau^2 \beta
\end{array}
\right],
\left[
\begin{array}{c}
\alpha \\ \eta_6 \\ \kappabar
\end{array}
\right]
\right\rangle
\]
is well-defined, where $\beta$ is an element of $\pi_{53,29}$
that is detected by $\D h_1 d_0^2$.
The same argument involving inclusion of the bottom cell into
$C\tau$ applies to this Toda bracket.

Finally, assume that $2 \alpha + \tau \eta \epsilon \eta_6$
is detected by $\tau^2 M d_0^2$.
Table \ref{tab:misc-extn} shows that $M d_0$ detects
$\kappa \theta_{4.5}$, so $\tau^2 M d_0^2$ detects 
$\tau^2 \kappa^2 \theta_{4.5}$.
Then $2 \alpha + \tau \eta \epsilon \eta_6$ equals either
$\tau^2 \kappa^2 \theta_{4.5}$ or 
$\tau^2 \kappa^2 \theta_{4.5} + \tau^2 \beta \kappabar$.
We can apply the same argument to the Toda bracket
\[
\left\langle
\eta,
\left[
\begin{array}{ccc}
2 & \tau \eta \epsilon & \tau^2 \kappa \theta_{4.5}
\end{array}
\right],
\left[
\begin{array}{c}
\alpha \\ \eta_6 \\ \kappa
\end{array}
\right]
\right\rangle,
\]
or to the Toda bracket
\[
\left\langle
\eta,
\left[
\begin{array}{cccc}
2 & \tau \eta \epsilon & \tau^2 \kappa \theta_{4.5} & \tau^2 \beta
\end{array}
\right],
\left[
\begin{array}{c}
\alpha \\ \eta_6 \\ \kappa \\ \kappabar
\end{array}
\right]
\right\rangle.
\]

We have now shown by contradiction that $h_4 D_2$ does not survive.
After ruling out other possibilities by comparison to $C\tau$
and to $\mmf$, the only remaining possibility is that 
$d_5(h_4 D_2)$ equals $\tau^4 d_1 g^2$.
\end{proof}

\begin{lemma}
\label{lem:d5-t^3gG0}
\revdeg{86, 11, 45}
$d_5(\tau^3 g G_0) = \tau M \D h_1^2 d_0$.
\end{lemma}

\begin{proof}
Suppose for sake of contradiction that the element
$\tau^3 g G_0$ survived.
It cannot be the target of a hidden $\tau$ extension,
and it cannot be hit by a differential.
Therefore, it maps non-trivially under inclusion of the bottom
cell into $C\tau$, and the only possible image is 
$\D^2 e_1 + \ol{\tau \D h_2 e_1 g}$.

Let $\alpha$ be an element of $\pi_{86,45}$ that is detected
by $\tau^3 g G_0$.  
Consider the Toda bracket
$\langle \alpha, 2\nu, \nu \rangle$.
Lemma \ref{lem:t^3gG0,h0h2,h2} implies that this Toda bracket
is detected by $e_0 x_{76,9}$, or is detected in higher Adams filtration.

On the other hand, 
under inclusion of the bottom cell into $C\tau$,
the Toda bracket is detected by $\D^2 h_1 g_2$.
This is inconsistent with the conclusion of the previous paragraph,
since inclusion of the bottom cell can only increase
Adams filtrations.

We now know that $\tau^3 g G_0$ does not survive.
After eliminating other possibilities by comparison to $\mmf$,
the only remaining possibility is that $d_5(\tau^3 g G_0)$ equals
$\tau M \D h_1^2 d_0$.
\end{proof}

\begin{lemma}
\revdeg{92, 4, 48}
\label{lem:d5-g3}
$d_5(g_3) = h_6 d_0^2$.
\end{lemma}

\begin{proof}
Table \ref{tab:Toda} shows that $h_1 h_6$ detects the Toda bracket
$\langle \eta, 2, \theta_5 \rangle$.  Therefore,
$h_1 h_6 d_0^2$ detects $\kappa^2 \langle \eta, 2, \theta_5 \rangle$.
Now consider the shuffle
\[
\tau \kappa^2 \langle \eta, 2, \theta_5 \rangle =
\langle \tau \kappa^2, \eta, 2 \rangle \theta_5.
\]
Lemma \ref{lem:tkappa^2,eta,2} shows that the last bracket is zero.
Therefore, $h_1 h_6 d_0^2$ does not support a hidden
$\tau$ extension, so it is either hit by a differential or 
in the image of projection from $C\tau$ to the top cell.

In the Adams spectral sequence for $C\tau$, the element 
$h_0^3 h_4^2 h_6$ detects the Toda bracket
$\langle \theta_4, 2, \theta_5 \rangle$.
Therefore, $h_0^3 h_4^2 h_6$ must be in the image of
inclusion of the bottom cell into $C\tau$.
In particular,
$h_0^3 h_4^2 h_6$ cannot map to $h_1 h_6 d_0^2$
under projection from $C\tau$ to the top cell.

Now $h_1 h_6 d_0^2$ cannot be in the image of
projection from $C\tau$ to the top cell, so it
must be hit by some differential.
The only possibility is that $d_5(h_1 g_3)$ equals $h_1 h_6 d_0^2$.
\end{proof}

\rev{
\begin{lemma}
\label{lem:d5-D^2g2}
\revdeg{92, 12, 48}
$d_5(\D^2 g_2) = 0$.
\end{lemma}

\begin{proof}
The only possible values for $d_2(\D^2 g_2)$ are the linear combinations
of $\tau M \D c_0 d_0 + \tau^2 M d_0 l$ 
and $\tau^2 \D^2 h_2 g^2$.
The possibilities
$\tau M \D c_0 d_0 + \tau^2 M d_0 l$ and
$\tau M \D c_0 d_0 + \tau^2 M d_0 l + \tau^2 \D^2 h_2 g^2$ are
ruled out by $d_0$ extensions.
More specifically,
$d_0 (\tau M \D c_0 d_0 + \tau^2 M d_0 l)$
and
$d_0(\tau M \D c_0 d_0 + \tau^2 M d_0 l + \tau^2 \D^2 h_2 g^2)$
equal the non-zero element $\tau^2 M d_0^2 l$
in the $E_5$-page, while
$d_0 \cdot \D^2 g_2$ is zero already in the $E_2$-page.

If $\tau^2 \D^2 h_2 g^2$ were the value of a differential,
then the $2$ extension from $\tau \D^2 h_2 g^2$ to $\tau \D^2 h_0 h_2 g^2$
would be detected by the top cell of $C\tau$.
However, there is no such $2$ extension in the homotopy of $C\tau$.
\end{proof}
}

\begin{lemma}
\label{lem:d5-e0x76,9}
\revdeg{93, 13, 50}
$d_5(e_0 x_{76,9}) = M \D h_1 c_0 d_0$.
\end{lemma}

\begin{proof}
If $M \D h_1 c_0 d_0$ were a permanent non-zero cycle, then it
could not support a hidden $\tau$ extension because
Lemma \ref{lem:hit-MPDh1d0} shows that $M P \D h_1 d_0$
is hit by some differential.
Therefore, it would lie in the image of projection from
$C\tau$ to the top cell, and the only possible pre-image is the
element $\D^2 h_1 g_2$ in the $E_\infty$-page of the
Adams spectral sequence for $C\tau$.

There is a $\sigma$ extension from $\D^2 e_1 + \ol{\tau \D h_2 e_1 g}$
to $\D^2 h_1 g_2$ in the Adams spectral sequence for $C\tau$.
Then $M \D h_1 c_0 d_0$ would also have to be the target of a 
$\sigma$ extension.
The only possible source for this extension would be $M \D h_1^2 d_0$.

Table \ref{tab:eta-extn} shows that $M h_1$
detects $\eta \theta_{4.5}$, so
$M \D h_1^2 d_0$ detects $\eta \theta_{4.5} \{\D h_1 d_0\}$.
The product $\eta \sigma \theta_{4.5} \{\D h_1 d_0\}$ equals zero
because $\sigma \{\D h_1 d_0\}$ is zero.
Therefore, $M \D h_1^2 d_0$ cannot support a hidden $\sigma$
extension to $M \D h_1 c_0 d_0$.

We have now shown that $M \D h_1 c_0 d_0$ must be hit by some
differential, and the only possibility is that equals 
$d_5(e_0 x_{76,9})$.
\end{proof}

\section{Higher differentials}
\label{sctn:Adams-higher}

\revv{
Table \ref{tab:Adams-higher} lists the multiplicative generators 
of the Adams $E_r$-page, for $r \geq 6$, through the 90-stem
whose $d_r$ differentials are non-zero, or whose $d_r$ differentials
are zero for non-obvious reasons.
}

\begin{thm}
\label{thm:Adams-higher}
Table \ref{tab:Adams-higher} lists some values of the
Adams $d_r$ differential on multiplicative generators of the
$E_r$-page, for $r \geq 6$.
For $r \geq 6$, 
the Adams $d_r$ differential is zero on all multiplicative
generators of the $E_r$-page not listed in the table.
The list is complete through the 90-stem, except that:
\rev{
\begin{enumerate}
\item
$d_{10}(h_1 f_2)$ might equal $M \D h_1 d_0$.
\item
$d_9(\tau x_{85,6} + h_0^3 c_3)$ might equal $\tau M \D h_1 d_0$.
\item
$d_9(h_4^2 D_3)$ might equal $\tau M \D h_1 g$.
\end{enumerate}
}
\end{thm}

\begin{proof}
The $d_r$ differential on many multiplicative generators is zero.
For the majority of such multiplicative generators,
the $d_r$ differential is zero
because there are no possible non-zero
values, or by comparison to the Adams spectral sequence for $C\tau$,
or by comparison to $\tmf$ or $\mmf$.
In a few cases, the multiplicative generator
is already known to be a permanent cycle, as shown in 
Table \ref{tab:Adams-perm}.
A few additional cases appear
in Table \ref{tab:Adams-higher} because their proofs require further
explanation.

Some of the more difficult computations appear in \cite{BIX}.
The remaining more difficult computations are carried out in the
following lemmas.
\end{proof}


\begin{lemma}
\label{lem:d6-tQ3+tn1}
\mbox{}
\begin{enumerate}
\item
\revdeg{67, 5, 35}
$d_6(\tau Q_3 + \tau n_1) = 0$.
\item
\revdeg{87, 9, 48}
$d_6(g Q_3) = 0$
\end{enumerate}
\end{lemma}

\begin{proof}
Several possible differentials on these elements are eliminated
by comparison to the
Adams spectral sequences for $C\tau$ and for $\tmf$.
The only remaining possibility is that $d_6(\tau Q_3 + \tau n_1)$
might equal $\tau^2 M h_1 g$, and
that $d_6(g Q_3)$ might equal $\tau M h_1 g^2$.

The element $M \D h_0^2 e_0$ is not hit by any differential because
Table \ref{tab:Adams-perm} shows that $h_1^2 c_3$ is a permanent cycle,
and Table \ref{tab:Toda} shows that $\tau^2 g Q_3 = h_4^2 Q_2$
must survive to detect the Toda bracket
$\langle \theta_4, \tau \kappabar, \{ t\} \rangle$.

Lemma \ref{lem:tetakappabar^2,2,4kappabar2}
shows that $M \D h_0^2 e_0$ 
detects the Toda bracket
$\langle \tau \eta \kappabar^2, 2, 4 \kappabar_2 \rangle$, which
contains
$\tau \kappabar^2 \langle \eta, 2, 4 \kappabar_2 \rangle$.
Lemma \ref{lem:eta,2,4kappabar2} shows that
this expression contains zero.
We now know that $M \D h_0^2 e_0$ detects an element
in the indeterminacy of the bracket
$\langle \tau \eta \kappabar^2, 4, 2 \kappabar_2 \rangle$.
In fact, it must detect a multiple
of $\tau \eta \kappabar^2$ since
$2 \kappabar_2 \cdot \pi_{42,22}$ is zero.

The only possibility is that $M \D h_0^2 e_0$
detects $\kappabar$ times an element detected by $\tau^2 M h_1 g$.
Therefore, $\tau^2 M h_1 g$ cannot be hit by a differential.
This shows that $\tau Q_3 + \tau n_1$ is a permanent cycle.

We also know that $M \D h_0^2 e_0$ is the target
of a hidden $\tau$ extension, since it detects a multiple of $\tau$.
The element $\tau^2 M h_1 g^2$
is the only possible source of this hidden $\tau$ extension,
so it cannot be hit by a differential.
This shows that $d_6(g Q_3)$ cannot equal $\tau M h_1 g^2$.
\end{proof}

\begin{lemma}
\label{lem:d6-h2^2H1}
\mbox{}
\begin{enumerate}
\item
\revdeg{68, 7, 37}
$d_6(h_2^2 H_1) = M c_0 d_0$.
\item
\revdeg{68, 7, 36}
$d_7(\tau h_2^2 H_1) = M P d_0$.
\end{enumerate}
\end{lemma}

\begin{proof}
Table \ref{tab:Massey} shows that 
$M P d_0$ equals the Massey product $\langle P d_0, h_0^3, g_2 \rangle$.
This implies that
$M P d_0$ detects the
Toda bracket $\langle \tau \eta^2 \kappabar, 8, \kappabar_2 \rangle$.
Lemma \ref{lem:teta^2kappabar,8,kappabar2} shows that this
Toda bracket consists entirely of multiples of
$\tau \eta^2 \kappabar$.

We now know that $M P d_0$ detects a multiple of 
$\tau \eta^2 \kappabar$.
The only possibility is that 
$M P d_0$ detects $\eta$ times an element detected by
$\tau^2 M h_1 g$.

We will show in Lemma \ref{lem:nu-th1H1}
that $\tau^2 M h_1 g$ is the target of a $\nu$ extension,
so $\tau^2 M h_1 g$ cannot support a hidden $\eta$ extension.
Therefore, $M P d_0$ must be hit by some differential.
The only possibility is that 
$d_7(\tau h_2^2 H_1)$ equals $M P d_0$.
Then $h_2^2 H_1$ cannot survive to the $E_7$-page,
so $d_6(h_2^2 H_1)$ equals $M c_0 d_0$.
\end{proof}

\begin{lemma}
\label{lem:perm-th1p1}
\revdeg{71, 5, 37}
The element $\tau h_1 p_1$ is a permanent cycle.
\end{lemma}

\begin{proof}
Lemma \ref{lem:d5-tp1+h0^2h3h6}, together with results of
\cite{BIX}, show that $\tau h_1 p_1$ survives to the $E_6$-page.
We must eliminate possible higher differentials.

Table \ref{tab:tau-extn} shows that there is a hidden
$\tau$ extension from $\tau h_2^2 C''$ to $\D^2 h_1^2 h_4 c_0$.
This means that $\tau h_2 C'' + h_1 h_3(\D e_1 + C_0)$ must 
also support a hidden $\tau$ extension.

The two possible targets for this hidden $\tau$ extension
are $\D^2 h_2 c_1$ and $\tau \D^2 h_1^2 g + \tau^3 \D h_2^2 g^2$.
The second possibility is ruled out by comparison to $\tmf$, so
$\D^2 h_2 c_1$ cannot be hit by a differential.
\end{proof}

\begin{lemma}
\label{lem:perm-Ph0h2h6}
\revdeg{74, 7, 38}
The element $P h_0 h_2 h_6$ is a permanent cycle.
\end{lemma}

\begin{proof}
First note that projection from $C\tau$ to the top cell
takes $P h_2 h_6$ to a non-zero element.
If $P h_0 h_2 h_6$ were not a permanent cycle in the
Adams spectral sequence for the sphere,
then projection from $C\tau$ to the top cell would
also take $P h_0 h_2 h_6$ to a non-zero element.
Then the $2$ extension from $P h_2 h_6$ to $P h_0 h_2 h_6$
in $\pi_{74,38} C\tau$ would project to a $2$ extension
in $\pi_{73,39}$.  However, there are no possible
$2$ extensions in $\pi_{73,39}$.
\end{proof}

\begin{lemma}
\label{lem:d7-m1}
\revdeg{77, 7, 42}
$d_7(m_1) = 0$.
\end{lemma}

\begin{proof}
The only other possibility is that $d_7(m_1)$ equals $\tau^2 g^2 t$.
If that were the case, then the $\nu$ extension from
$\tau g^2 t$ to $\tau^2 c_1 g^3$ would be detected by projection
from $C\tau$ to the top cell.  However, the homotopy groups of
$C\tau$ have no such $\nu$ extension.
\end{proof}

\begin{lemma}
\label{lem:d8-h1x1}
\revdeg{80, 6, 43}
$d_8(h_1 x_1) = 0$.
\end{lemma}

\begin{proof}
Table \ref{tab:Adams-perm} shows that $\tau h_1 x_1$ is a permanent
cycle.
Then $d_8(\tau h_1 x_1)$ cannot equal 
$\tau^2 M e_0^2$, and $d_8(h_1 x_1)$ cannot equal
$\tau M e_0^2$.
\end{proof}

\begin{lemma}
\label{lem:d6-h2h4h6}
\revdeg{81, 3, 42}
$d_6(h_2 h_4 h_6) = 0$.
\end{lemma}

\begin{proof}
Table \ref{tab:Adams-perm} shows that $h_2^2 h_4 h_6$ is a permanent
cycle.  Therefore, the Adams differential
$d_6(h_2^2 h_4 h_6)$ does not equal
$\tau h_2 c_1 A'$, and $d_6(h_2 h_4 h_6)$ does not equal
$\tau c_1 A'$.
\end{proof}

\rev{
\begin{lemma}
\label{lem:d10-h2h6g+h1^2f2}
\revdeg{86, 6, 46}
$d_{10}(h_2 h_6 g + h_1^2 f_2) = 0$.
\end{lemma}

\begin{proof}
If $d_{10}(h_2 h_6 g + h_1^2 f_2)$ equaled
$M \D h_1^2 d_0$, then the $\nu$ extension from
$\tau h_6 g + \tau h_2 e_2$ to $\tau h_2 h_6 g + \tau h_1^2 f_2$
would be detected by the bottom cell of $C\tau$.
However, there is no such $\nu$ extension in the homotopy 
of $C\tau$.
\end{proof}
}


\begin{lemma}
\label{lem:d7-x87,7}
\revdeg{87, 7, 45}
$d_7(x_{87,7}) = 0$.
\end{lemma}

\begin{proof}
If $\tau \D^2 h_2^2 d_1$ were hit by a differential, then
the $\nu$ extension from $\D^2 h_2^2 d_1$ to $\D^2 h_1^2 h_3 d_1$
would be detected by projection from $C\tau$ to the top cell.
But the homotopy of $C\tau$ has no such $\nu$ extension.
\end{proof}

\rev{
\begin{lemma}
\label{lem:d6-tDh1H1}
\revdeg{87, 10, 45}
$d_6(\tau \D h_1 H_1) = \tau M \D h_0^2 e_0$.
\end{lemma}

\begin{proof}
Suppose for sake of contradiction that
$\tau \D h_1 H_1$ survives.
This element cannot be hit, nor can it be the target of a hidden
$\tau$ extension.  Therefore, it would have non-zero image 
under inclusion of the bottom cell into $C\tau$, and it would
map to $\ol{\D h_1 B_7}$.

In the homotopy of $C\tau$, the Toda bracket
$\langle \ol{\D h_1 B_7}, h_0, h_2^2 \rangle$
is detected by the element $M \D^2 h_1$.
Beware that this bracket has indeterminacy in lower filtration
since $h_3 \cdot \ol{\D h_1 B_7} = \ol{h_1^2 x_{91,11}}$.

This implies that the Toda bracket
$\langle \{\tau \D h_1 H_1\}, 2, \nu^2 \rangle$
would be non-zero in $\pi_{94,49}$, and all of its elements 
would be detected in Adams filtration at most $15$.
(Beware that this Toda bracket would have indeterminacy
detected by $\tau \D h_1 h_3 H_1$.)

On the other hand, the Moss Convergence Theorem \ref{thm:Moss}
would imply that the Toda bracket is detected in filtration
at least $12$.
However, there are no possible elements in filtrations $12$ through $15$.

We have now shown that $\tau \D h_1 H_1$ cannot survive.
There is only one possible value for a differential on
$\tau \D h_1 H_1$.

The previous argument assumed that $2 \{\tau \D h_1 H_1\}$ is zero
in order to form the Toda bracket
$\langle \{\tau \D h_1 H_1\}, 2, \nu^2 \rangle$.
However, it is possible that $\tau \D h_1 H_1$ supports a hidden
$2$ extension to $\tau \D^2 h_3 d_1$ or to $\tau^2 \D^2 c_1 g$.
Therefore, $2 \{\tau \D h_1 H_1\}$ might equal
$\tau \sigma \{\D^2 d_1\}$, $\tau \nu \{\D^2 t\}$, or their sum.
In those cases, we would need to consider the matric Toda brackets
\[
\left\langle 
\left[
\begin{array}{cc}
\{ \tau \D h_1 H_1\} & \tau \{\D^2 d_1\}
\end{array}
\right],
\left[
\begin{array}{c}
2 \\ \sigma
\end{array}
\right],
\nu^2
\right\rangle,
\]
\[
\left\langle 
\left[
\begin{array}{cc}
\{ \tau \D h_1 H_1\} & \nu
\end{array}
\right],
\left[
\begin{array}{c}
2 \\ \tau \{\D^2 t\}
\end{array}
\right],
\nu^2
\right\rangle,
\]
or
\[
\left\langle 
\left[
\begin{array}{ccc}
\{ \tau \D h_1 H_1\} & \tau \{\D^2 d_1\} & \nu
\end{array}
\right],
\left[
\begin{array}{c}
2 \\ \sigma \\ \tau \{\D^2 t \}
\end{array}
\right],
\nu^2
\right\rangle
\]
respectively.
Under inclusion of the bottom cell into $C\tau$, these three brackets
would map to
\[
\left\langle 
\left[
\begin{array}{cc}
\ol{\D h_1 B_7} & 0
\end{array}
\right],
\left[
\begin{array}{c}
h_0 \\ h_3
\end{array}
\right],
h_2^2
\right\rangle,
\]
\[
\left\langle 
\left[
\begin{array}{cc}
\ol{\D h_1 B_7} & h_2
\end{array}
\right],
\left[
\begin{array}{c}
h_0 \\ 0
\end{array}
\right],
h_2^2
\right\rangle,
\]
or
\[
\left\langle 
\left[
\begin{array}{ccc}
\ol{\D h_1 B_7} & 0 & h_2
\end{array}
\right],
\left[
\begin{array}{c}
h_0 \\ h_3 \\ 0
\end{array}
\right],
h_2^2
\right\rangle
\]
respectively.
All three of these brackets in $C\tau$ equal
$\langle \ol{\D h_1 B_7}, h_0, h_2^2 \rangle$.
Beware that the last two could have larger indeterminacy,
but in fact do not.
\end{proof}
}

\begin{lemma}
\label{lem:perm-x88,10}
\revdeg{88, 10, 48}
The element $x_{88,10}$ is a permanent cycle.
\end{lemma}

\begin{proof}
In the Adams spectral sequence for $C\tau$, there is a hidden
$\eta$ extension from $h_1^2 x_{85,6}$ to $x_{88,10}$.
Therefore, $x_{88,10}$ lies in the image of inclusion of the
bottom cell into $C\tau$.
The only possible pre-image is the element $x_{88,10}$ in the
Adams spectral sequence in the sphere, so $x_{88,10}$ must survive.
\end{proof}

\begin{lemma}
\label{lem:d6-h2^2gH1}
\mbox{}
\begin{enumerate}
\item
\revdeg{88, 11, 49}
$d_6(h_2^2 g H_1) = M c_0 e_0^2$.
\item
\revdeg{88, 11, 48}
$d_7(\tau h_2^2 g H_1) = 0$.
\end{enumerate}
\end{lemma}

\begin{proof}
If $M c_0 e_0^2$ is non-zero in the $E_\infty$-page, then it detects
an element that is annihilated by $\tau$ because Lemma \ref{lem:d6-Dg2g}
shows that the only possible target of such an extension is hit
by a differential.
Then $M c_0 e_0^2$ would be in the image of projection 
from $C\tau$ to the top cell.
The only possible pre-image would be the element $\D g_2 g$ of the Adams
spectral sequence for $C\tau$.

In the Adams spectral sequence for $C\tau$, there is a $\sigma$
extension from $g A'$ to $\D g_2 g$.  
Projection from $C\tau$ to the top cell would imply that there
is a hidden $\sigma$ extension in the homotopy groups of the sphere,
from $M h_1 e_0^2$ to $M c_0 e_0^2$, because
$g A'$ maps to $M h_1 e_0^2$ under projection from $C\tau$
to the top cell.

But $M h_1 e_0^2$ detects $\eta \theta_{4.5} \{ e_0^2 \}$,
which cannot support a $\sigma$ extension.  This establishes the
first formula.

For the second formula, 
if $d_7(\tau h_2^2 g H_1)$ were equal to $\tau^2 \D h_2^2 e_0 g^2$,
then the same argument would apply,
with $\tau \D h_2^2 e_0 g^2$ substituted for $M c_0 e_0^2$.
\end{proof}

\begin{lemma}
\label{lem:d6-Dg2g}
\revdeg{88, 12, 48}
$d_6(\D g_2 g) = M d_0^3$.
\end{lemma}

\begin{proof}
The proof of Lemma \ref{lem:d4-th2B5g} shows that
$M d_0^3$ must be hit by a differential.
The only possibility is that $d_6(\D g_2 g)$ equals
$M d_0^3$.

Alternatively, Lemma \ref{lem:d6-h2^2H1} shows that
$d_7(\tau h_2^2 H_1) = M P d_0$.
Note that $\tau g \cdot \tau h_2^2 H_1 = 0$ in the $E_7$-page.
Therefore, $\tau M d_0^3 = \tau g \cdot M P d_0$ must already
be zero in the $E_7$-page.  The only possibility is that
$d_6(\tau \D g_2 g) = \tau M d_0^3$, and then
$d_6(\D g_2 g) = M d_0^3$.
\end{proof}

\rev{
\begin{remark}
\label{rem:d6-D^2f1}
Table \ref{tab:Adams-higher} shows that
$d_6(\D^2 f_1)$ equals $\tau^2 M d_0^3$.
The proof relies on $d_6(\tau \D h_1 H_1) = \tau M \D h_0^2 e_0$
and uses techniques similar to the ones in \cite{Chua21}.
\end{remark}
}

\begin{lemma}
\label{lem:perm-h0g3}
\revdeg{92, 5, 48}
The element $h_0 g_3$ is a permanent cycle.
\end{lemma}

\begin{proof}
In the homotopy of $C\tau$, the product $\theta_4 \theta_5$ is detected
by $h_0^2 g_3$.  In the sphere, the product $\theta_4 \theta_5$ 
is therefore non-zero and detected in Adams filtration at most 6.

Table \ref{tab:Toda} shows that 
the Toda bracket $\langle 2, \theta_4, \theta_4, 2 \rangle$
contains $\theta_5$.
Therefore, the product $\theta_4 \theta_5$ is contained in
\[
\theta_4 \langle 2, \theta_4, \theta_4, 2 \rangle =
\langle \theta_4, 2, \theta_4, \theta_4 \rangle 2.
\]
(Note that the sub-bracket $\langle \theta_4, \theta_4, 2 \rangle$
is zero because $\pi_{61,32}$ is zero.)
Therefore, $\theta_4 \theta_5$ is divisible by $2$.
It follows that
$\theta_4 \theta_5$ is detected by $h_0^2 g_3$, and
$h_0 g_3$ is a permanent cycle that detects
$\langle \theta_4, 2, \theta_4, \theta_4 \rangle$.
\end{proof}


\begin{lemma}
\label{lem:d6-D1h1^2e1}
\revdeg{92, 10, 51}
$d_6(\D_1 h_1^2 e_1) = 0$.
\end{lemma}

\begin{proof}
Consider the element $\ol{\tau M h_2^2 g^2}$
in the Adams spectral sequence for $C\tau$.
This element cannot be in the image of inclusion of the bottom
cell into $C\tau$.  Therefore,
it must map non-trivially under projection from $C\tau$
to the top cell.
The only possibility is that
$\tau M h_2^2 g^2$ is the image.
Therefore, $\tau M h_2^2 g^2$ cannot be the target of a differential.
\end{proof}

\begin{lemma}
\label{lem:d7-x92,10}
\revdeg{92, 10, 48}
$d_7(x_{92,10})$ does not equal $\tau^2 \D^2 h_2 g^2$.
\end{lemma}

\begin{proof}
If $\tau^2 \D^2 h_2 g^2$ were hit by a differential, then
the $2$ extension from $\tau \D^2 h_2 g^2$ to $\tau \D^2 h_0 h_2 g^2$
would be detected by projection from $C\tau$ to the top cell.
But the homotopy of $C\tau$ has no such $2$ extension.
\end{proof}

\begin{lemma}
\label{lem:d8-D1h2e1}
\revdeg{92, 10, 51}
$d_8(\D_1 h_2 e_1) = 0$.
\end{lemma}

\begin{proof}
Consider the element $e_0 x_{76,9}$ in the
Adams $E_\infty$-page for $C\tau$.
It cannot be in the image of inclusion of the bottom cell
into $C\tau$, so it must project to a non-zero element in the
top cell.  The only possible image is $M \D h_1^3 g$.
Therefore, $M \D h_1^3 g$ cannot be the target of a differential.
\end{proof}

\begin{lemma}
\label{lem:hit-MPDh1d0}
The element $M P \D h_1 d_0$ is hit by some differential.
\end{lemma}

\begin{proof}
Table \ref{tab:tau-extn} shows that
there is a hidden $\tau$ extension from
$\D h_1 c_0 d_0$ to $P \D h_1 d_0$.  Therefore,
$P \D h_1 d_0$ detects $\tau \epsilon \{\D h_1 d_0\}$.
On the other hand, Tables \ref{tab:eta-extn} and
\ref{tab:misc-extn} show that
$P \D h_1 d_0$ also detects $\tau \eta \kappa \{\D h_1 h_3\}$.
Since there are no elements in higher Adams filtration, we 
have that $\tau \epsilon \{\D h_1 d_0\}$ equals
$\tau \eta \kappa \{\D h_1 h_3\}$.

Table \ref{tab:misc-extn} shows that
$M P$ detects $\tau \epsilon \theta_{4.5}$,
so $M P \D h_1 d_0$ detects 
$\tau \epsilon \{\D h_1 d_0\} \theta_{4.5}$, 
which equals 
$\tau \eta \kappa \{\D h_1 h_3\} \theta_{4.5}$.
But $\tau \eta \kappa \theta_{4.5}$ is zero because
all elements of $\pi_{60,32}$ are detected by $\tmf$.
This shows that $M P \D h_1 d_0$ detects zero, so it
must be hit by a differential.
\end{proof}

\rev{
\begin{remark}
\label{rem:d6-te0x76,9}
Lemma \ref{lem:hit-MPDh1d0} does not specify the differential that
hits the element $M P \D h_1 d_0$.  In fact, $d_6(\tau e_0 x_{76,9})$ equals
$M P \D h_1 d_0$ \cite{BIX}.
\end{remark}
}